% Options for packages loaded elsewhere
\PassOptionsToPackage{unicode}{hyperref}
\PassOptionsToPackage{hyphens}{url}
\documentclass[
]{article}
\usepackage{xcolor}
\usepackage[margin=1in]{geometry}
\usepackage{amsmath,amssymb}
\setcounter{secnumdepth}{-\maxdimen} % remove section numbering
\usepackage{iftex}
\ifPDFTeX
  \usepackage[T1]{fontenc}
  \usepackage[utf8]{inputenc}
  \usepackage{textcomp} % provide euro and other symbols
\else % if luatex or xetex
  \usepackage{unicode-math} % this also loads fontspec
  \defaultfontfeatures{Scale=MatchLowercase}
  \defaultfontfeatures[\rmfamily]{Ligatures=TeX,Scale=1}
\fi
\usepackage{lmodern}
\ifPDFTeX\else
  % xetex/luatex font selection
\fi
% Use upquote if available, for straight quotes in verbatim environments
\IfFileExists{upquote.sty}{\usepackage{upquote}}{}
\IfFileExists{microtype.sty}{% use microtype if available
  \usepackage[]{microtype}
  \UseMicrotypeSet[protrusion]{basicmath} % disable protrusion for tt fonts
}{}
\makeatletter
\@ifundefined{KOMAClassName}{% if non-KOMA class
  \IfFileExists{parskip.sty}{%
    \usepackage{parskip}
  }{% else
    \setlength{\parindent}{0pt}
    \setlength{\parskip}{6pt plus 2pt minus 1pt}}
}{% if KOMA class
  \KOMAoptions{parskip=half}}
\makeatother
\usepackage{color}
\usepackage{fancyvrb}
\newcommand{\VerbBar}{|}
\newcommand{\VERB}{\Verb[commandchars=\\\{\}]}
\DefineVerbatimEnvironment{Highlighting}{Verbatim}{commandchars=\\\{\}}
% Add ',fontsize=\small' for more characters per line
\usepackage{framed}
\definecolor{shadecolor}{RGB}{248,248,248}
\newenvironment{Shaded}{\begin{snugshade}}{\end{snugshade}}
\newcommand{\AlertTok}[1]{\textcolor[rgb]{0.94,0.16,0.16}{#1}}
\newcommand{\AnnotationTok}[1]{\textcolor[rgb]{0.56,0.35,0.01}{\textbf{\textit{#1}}}}
\newcommand{\AttributeTok}[1]{\textcolor[rgb]{0.13,0.29,0.53}{#1}}
\newcommand{\BaseNTok}[1]{\textcolor[rgb]{0.00,0.00,0.81}{#1}}
\newcommand{\BuiltInTok}[1]{#1}
\newcommand{\CharTok}[1]{\textcolor[rgb]{0.31,0.60,0.02}{#1}}
\newcommand{\CommentTok}[1]{\textcolor[rgb]{0.56,0.35,0.01}{\textit{#1}}}
\newcommand{\CommentVarTok}[1]{\textcolor[rgb]{0.56,0.35,0.01}{\textbf{\textit{#1}}}}
\newcommand{\ConstantTok}[1]{\textcolor[rgb]{0.56,0.35,0.01}{#1}}
\newcommand{\ControlFlowTok}[1]{\textcolor[rgb]{0.13,0.29,0.53}{\textbf{#1}}}
\newcommand{\DataTypeTok}[1]{\textcolor[rgb]{0.13,0.29,0.53}{#1}}
\newcommand{\DecValTok}[1]{\textcolor[rgb]{0.00,0.00,0.81}{#1}}
\newcommand{\DocumentationTok}[1]{\textcolor[rgb]{0.56,0.35,0.01}{\textbf{\textit{#1}}}}
\newcommand{\ErrorTok}[1]{\textcolor[rgb]{0.64,0.00,0.00}{\textbf{#1}}}
\newcommand{\ExtensionTok}[1]{#1}
\newcommand{\FloatTok}[1]{\textcolor[rgb]{0.00,0.00,0.81}{#1}}
\newcommand{\FunctionTok}[1]{\textcolor[rgb]{0.13,0.29,0.53}{\textbf{#1}}}
\newcommand{\ImportTok}[1]{#1}
\newcommand{\InformationTok}[1]{\textcolor[rgb]{0.56,0.35,0.01}{\textbf{\textit{#1}}}}
\newcommand{\KeywordTok}[1]{\textcolor[rgb]{0.13,0.29,0.53}{\textbf{#1}}}
\newcommand{\NormalTok}[1]{#1}
\newcommand{\OperatorTok}[1]{\textcolor[rgb]{0.81,0.36,0.00}{\textbf{#1}}}
\newcommand{\OtherTok}[1]{\textcolor[rgb]{0.56,0.35,0.01}{#1}}
\newcommand{\PreprocessorTok}[1]{\textcolor[rgb]{0.56,0.35,0.01}{\textit{#1}}}
\newcommand{\RegionMarkerTok}[1]{#1}
\newcommand{\SpecialCharTok}[1]{\textcolor[rgb]{0.81,0.36,0.00}{\textbf{#1}}}
\newcommand{\SpecialStringTok}[1]{\textcolor[rgb]{0.31,0.60,0.02}{#1}}
\newcommand{\StringTok}[1]{\textcolor[rgb]{0.31,0.60,0.02}{#1}}
\newcommand{\VariableTok}[1]{\textcolor[rgb]{0.00,0.00,0.00}{#1}}
\newcommand{\VerbatimStringTok}[1]{\textcolor[rgb]{0.31,0.60,0.02}{#1}}
\newcommand{\WarningTok}[1]{\textcolor[rgb]{0.56,0.35,0.01}{\textbf{\textit{#1}}}}
\usepackage{longtable,booktabs,array}
\newcounter{none} % for unnumbered tables
\usepackage{calc} % for calculating minipage widths
% Correct order of tables after \paragraph or \subparagraph
\usepackage{etoolbox}
\makeatletter
\patchcmd\longtable{\par}{\if@noskipsec\mbox{}\fi\par}{}{}
\makeatother
% Allow footnotes in longtable head/foot
\IfFileExists{footnotehyper.sty}{\usepackage{footnotehyper}}{\usepackage{footnote}}
\makesavenoteenv{longtable}
\usepackage{graphicx}
\makeatletter
\newsavebox\pandoc@box
\newcommand*\pandocbounded[1]{% scales image to fit in text height/width
  \sbox\pandoc@box{#1}%
  \Gscale@div\@tempa{\textheight}{\dimexpr\ht\pandoc@box+\dp\pandoc@box\relax}%
  \Gscale@div\@tempb{\linewidth}{\wd\pandoc@box}%
  \ifdim\@tempb\p@<\@tempa\p@\let\@tempa\@tempb\fi% select the smaller of both
  \ifdim\@tempa\p@<\p@\scalebox{\@tempa}{\usebox\pandoc@box}%
  \else\usebox{\pandoc@box}%
  \fi%
}
% Set default figure placement to htbp
\def\fps@figure{htbp}
\makeatother
\setlength{\emergencystretch}{3em} % prevent overfull lines
\providecommand{\tightlist}{%
  \setlength{\itemsep}{0pt}\setlength{\parskip}{0pt}}
\usepackage{bookmark}
\IfFileExists{xurl.sty}{\usepackage{xurl}}{} % add URL line breaks if available
\urlstyle{same}
\hypersetup{
  pdftitle={Projekt - Airline Loyalty Program Promotion and Beyond},
  pdfauthor={Manca Prošek, Tristan Dragan, Patricija Jaklič},
  hidelinks,
  pdfcreator={LaTeX via pandoc}}

\title{Projekt - Airline Loyalty Program Promotion and Beyond}
\author{Manca Prošek, Tristan Dragan, Patricija Jaklič}
\date{2026-08-14}

\begin{document}
\maketitle

\section{Branje podatkov}\label{branje-podatkov}

\subsection{1. Customer Flight
Activity.csv}\label{customer-flight-activity.csv}

Osnovne lastnosti podatkov (Customer Flight Activity): Št. vrstic:
392,936 Št. stolpcev: 8 Ni manjkajočih podatkov/vrstic Vse spremenljivke
so numerične Vse spremenljivke razen Loyalty Number in Month so
asimetrične v desno.

\begin{Shaded}
\begin{Highlighting}[]
\NormalTok{pregled\_CFA }\OtherTok{\textless{}{-}} \FunctionTok{create\_report}\NormalTok{(Customer\_Flight\_Activity)}
\end{Highlighting}
\end{Shaded}

\begin{verbatim}
## 
## 
## processing file: report.rmd
\end{verbatim}

\begin{verbatim}
##   |                                             |                                     |   0%  |                                             |.                                    |   2%                                   |                                             |..                                   |   5% [global_options]                  |                                             |...                                  |   7%                                   |                                             |....                                 |  10% [introduce]                       |                                             |....                                 |  12%                                   |                                             |.....                                |  14% [plot_intro]
\end{verbatim}

\begin{verbatim}
##   |                                             |......                               |  17%                                   |                                             |.......                              |  19% [data_structure]                  |                                             |........                             |  21%                                   |                                             |.........                            |  24% [missing_profile]
\end{verbatim}

\begin{verbatim}
##   |                                             |..........                           |  26%                                   |                                             |...........                          |  29% [univariate_distribution_header]  |                                             |...........                          |  31%                                   |                                             |............                         |  33% [plot_histogram]
\end{verbatim}

\begin{verbatim}
##   |                                             |.............                        |  36%                                   |                                             |..............                       |  38% [plot_density]                    |                                             |...............                      |  40%                                   |                                             |................                     |  43% [plot_frequency_bar]              |                                             |.................                    |  45%                                   |                                             |..................                   |  48% [plot_response_bar]               |                                             |..................                   |  50%                                   |                                             |...................                  |  52% [plot_with_bar]                   |                                             |....................                 |  55%                                   |                                             |.....................                |  57% [plot_normal_qq]
\end{verbatim}

\begin{verbatim}
##   |                                             |......................               |  60%                                   |                                             |.......................              |  62% [plot_response_qq]                |                                             |........................             |  64%                                   |                                             |.........................            |  67% [plot_by_qq]                      |                                             |..........................           |  69%                                   |                                             |..........................           |  71% [correlation_analysis]
\end{verbatim}

\begin{verbatim}
##   |                                             |...........................          |  74%                                   |                                             |............................         |  76% [principal_component_analysis]
\end{verbatim}

\begin{verbatim}
##   |                                             |.............................        |  79%                                   |                                             |..............................       |  81% [bivariate_distribution_header]   |                                             |...............................      |  83%                                   |                                             |................................     |  86% [plot_response_boxplot]           |                                             |.................................    |  88%                                   |                                             |.................................    |  90% [plot_by_boxplot]                 |                                             |..................................   |  93%                                   |                                             |...................................  |  95% [plot_response_scatterplot]       |                                             |.................................... |  98%                                   |                                             |.....................................| 100% [plot_by_scatterplot]           
\end{verbatim}

\begin{verbatim}
## output file: /Users/patricijajaklic/Documents/FAKS - statistika/2025:2026/Poslovna statistika/Poslovna statistik- Airline/report.knit.md
\end{verbatim}

\begin{verbatim}
## /private/var/folders/51/r10wy95d4vx03_k0ln6yflrr0000gn/T/AppTranslocation/4815461A-3C6C-416C-999D-85E63FCF6DA7/d/RStudio.app/Contents/Resources/app/quarto/bin/tools/aarch64/pandoc +RTS -K512m -RTS '/Users/patricijajaklic/Documents/FAKS - statistika/2025:2026/Poslovna statistika/Poslovna statistik- Airline/report.knit.md' --to html4 --from markdown+autolink_bare_uris+tex_math_single_backslash --output pandocb3b16e4d4c9e.html --lua-filter /Library/Frameworks/R.framework/Versions/4.5-arm64/Resources/library/rmarkdown/rmarkdown/lua/pagebreak.lua --lua-filter /Library/Frameworks/R.framework/Versions/4.5-arm64/Resources/library/rmarkdown/rmarkdown/lua/latex-div.lua --lua-filter /Library/Frameworks/R.framework/Versions/4.5-arm64/Resources/library/rmarkdown/rmarkdown/lua/table-classes.lua --embed-resources --standalone --variable bs3=TRUE --section-divs --table-of-contents --toc-depth 6 --template /Library/Frameworks/R.framework/Versions/4.5-arm64/Resources/library/rmarkdown/rmd/h/default.html --syntax-highlighting none --variable highlightjs=1 --variable theme=yeti --mathjax --variable 'mathjax-url=https://mathjax.rstudio.com/latest/MathJax.js?config=TeX-AMS-MML_HTMLorMML' --include-in-header /var/folders/51/r10wy95d4vx03_k0ln6yflrr0000gn/T//Rtmp0hECjK/rmarkdown-strb3b116bb7a1a.html
\end{verbatim}

\begin{verbatim}
## 
## Output created: report.html
\end{verbatim}

\subsection{2. Customer Loyalty
History.csv}\label{customer-loyalty-history.csv}

Osnovne lastnosti podatkov (Customer Loyalty History): Št. vrstic:
16,737 Št. stolpcev: 16 Št. manjkajočih podatkov: 33,578 Št. popolnih
vrstic: 1,540 Št. numeričnih spremenljivk: 7 Št. opisnih spremenljivk: 9
Spremenljivke z manjkajočimi vrednostmi so: Salary (25.32\%),
Cancellation Month (87.65\%) in Cancellation Year (87.65\%).

\begin{Shaded}
\begin{Highlighting}[]
\NormalTok{pregled\_CLH }\OtherTok{\textless{}{-}} \FunctionTok{create\_report}\NormalTok{(Customer\_Loyalty\_History)}
\end{Highlighting}
\end{Shaded}

\begin{verbatim}
## 
## 
## processing file: report.rmd
\end{verbatim}

\begin{verbatim}
##   |                                             |                                     |   0%  |                                             |.                                    |   2%                                   |                                             |..                                   |   5% [global_options]                  |                                             |...                                  |   7%                                   |                                             |....                                 |  10% [introduce]                       |                                             |....                                 |  12%                                   |                                             |.....                                |  14% [plot_intro]
\end{verbatim}

\begin{verbatim}
##   |                                             |......                               |  17%                                   |                                             |.......                              |  19% [data_structure]                  |                                             |........                             |  21%                                   |                                             |.........                            |  24% [missing_profile]
\end{verbatim}

\begin{verbatim}
##   |                                             |..........                           |  26%                                   |                                             |...........                          |  29% [univariate_distribution_header]  |                                             |...........                          |  31%                                   |                                             |............                         |  33% [plot_histogram]
\end{verbatim}

\begin{verbatim}
##   |                                             |.............                        |  36%                                   |                                             |..............                       |  38% [plot_density]                    |                                             |...............                      |  40%                                   |                                             |................                     |  43% [plot_frequency_bar]
\end{verbatim}

\begin{verbatim}
##   |                                             |.................                    |  45%                                   |                                             |..................                   |  48% [plot_response_bar]               |                                             |..................                   |  50%                                   |                                             |...................                  |  52% [plot_with_bar]                   |                                             |....................                 |  55%                                   |                                             |.....................                |  57% [plot_normal_qq]
\end{verbatim}

\begin{verbatim}
##   |                                             |......................               |  60%                                   |                                             |.......................              |  62% [plot_response_qq]                |                                             |........................             |  64%                                   |                                             |.........................            |  67% [plot_by_qq]                      |                                             |..........................           |  69%                                   |                                             |..........................           |  71% [correlation_analysis]
\end{verbatim}

\begin{verbatim}
##   |                                             |...........................          |  74%                                   |                                             |............................         |  76% [principal_component_analysis]
\end{verbatim}

\begin{verbatim}
##   |                                             |.............................        |  79%                                   |                                             |..............................       |  81% [bivariate_distribution_header]   |                                             |...............................      |  83%                                   |                                             |................................     |  86% [plot_response_boxplot]           |                                             |.................................    |  88%                                   |                                             |.................................    |  90% [plot_by_boxplot]                 |                                             |..................................   |  93%                                   |                                             |...................................  |  95% [plot_response_scatterplot]       |                                             |.................................... |  98%                                   |                                             |.....................................| 100% [plot_by_scatterplot]           
\end{verbatim}

\begin{verbatim}
## output file: /Users/patricijajaklic/Documents/FAKS - statistika/2025:2026/Poslovna statistika/Poslovna statistik- Airline/report.knit.md
\end{verbatim}

\begin{verbatim}
## /private/var/folders/51/r10wy95d4vx03_k0ln6yflrr0000gn/T/AppTranslocation/4815461A-3C6C-416C-999D-85E63FCF6DA7/d/RStudio.app/Contents/Resources/app/quarto/bin/tools/aarch64/pandoc +RTS -K512m -RTS '/Users/patricijajaklic/Documents/FAKS - statistika/2025:2026/Poslovna statistika/Poslovna statistik- Airline/report.knit.md' --to html4 --from markdown+autolink_bare_uris+tex_math_single_backslash --output pandocb3b15465495e.html --lua-filter /Library/Frameworks/R.framework/Versions/4.5-arm64/Resources/library/rmarkdown/rmarkdown/lua/pagebreak.lua --lua-filter /Library/Frameworks/R.framework/Versions/4.5-arm64/Resources/library/rmarkdown/rmarkdown/lua/latex-div.lua --lua-filter /Library/Frameworks/R.framework/Versions/4.5-arm64/Resources/library/rmarkdown/rmarkdown/lua/table-classes.lua --embed-resources --standalone --variable bs3=TRUE --section-divs --table-of-contents --toc-depth 6 --template /Library/Frameworks/R.framework/Versions/4.5-arm64/Resources/library/rmarkdown/rmd/h/default.html --syntax-highlighting none --variable highlightjs=1 --variable theme=yeti --mathjax --variable 'mathjax-url=https://mathjax.rstudio.com/latest/MathJax.js?config=TeX-AMS-MML_HTMLorMML' --include-in-header /var/folders/51/r10wy95d4vx03_k0ln6yflrr0000gn/T//Rtmp0hECjK/rmarkdown-strb3b13c4cc627.html
\end{verbatim}

\begin{verbatim}
## 
## Output created: report.html
\end{verbatim}

\section{Viri}\label{viri}

\subsection{Customer
churn-----------------------------------------}\label{customer-churn}

\subsection{1. članek/vir (osnove customer churn): Customer Churn
Prediction: A Systematic Review of Recent Advances, Trends, and
Challenges in Machine Learning and Deep
Learning}\label{ux10dlanekvir-osnove-customer-churn-customer-churn-prediction-a-systematic-review-of-recent-advances-trends-and-challenges-in-machine-learning-and-deep-learning}

\emph{Imani, M., Joudaki, M., Beikmohammadi, A., \& Arabnia, H. R.
(2025)} Link: {[}1.clanek{]}
(www.mdpi.com/2504-4990/7/3/105?utm\_source=chatgpt.com)

\subsection{2. članek/vir: An asymmetrical approach to understanding con
gurations of customer loyalty in the airline
industry}\label{ux10dlanekvir-an-asymmetrical-approach-to-understanding-con-gurations-of-customer-loyalty-in-the-airline-industry}

\emph{Catherine Prenticea, Sandra Maria Correia Loureirob} Link:
{[}2.clanek{]}
(\url{https://www.sciencedirect.com/science/article/pii/S0969698917301789?fr=RR-1&ref=cra_js_challenge})

\subsection{3. članek/vir: Key drivers of passenger loyalty: A case of
Frankfurt--Istanbul flights.
(2016).}\label{ux10dlanekvir-key-drivers-of-passenger-loyalty-a-case-of-frankfurtistanbul-flights.-2016.}

\emph{Journal of Air Transport Management} Link: {[}3.clanek{]}
(\url{https://www.sciencedirect.com/science/article/pii/S0969699715300120?utm_source=chatgpt.com})

\subsection{4. članek/vir: Frequent flyer programs in academic studies:
A literature review of research progress since 2000.
(2025).}\label{ux10dlanekvir-frequent-flyer-programs-in-academic-studies-a-literature-review-of-research-progress-since-2000.-2025.}

\emph{Journal of Air Transport Management.} Link: {[}4.clanek{]}
(\url{https://www.sciencedirect.com/science/article/pii/S0969699725001838?utm_source=chatgpt.com})

\subsection{5. članek/vir: In Pursuit of Enhanced Customer Retention
Management: Review, Key Issues, and Future
Directions}\label{ux10dlanekvir-in-pursuit-of-enhanced-customer-retention-management-review-key-issues-and-future-directions}

\emph{Ascarza, E., Neslin, S. A., Netzer, O., et al.~(2018)} Link:
{[}5.clanek{]}
(\url{https://business.columbia.edu/sites/default/files-efs/pubfiles/25780/ascarza_et_al_choice_symposium.pdf})

\subsection{6. članek/vir: Beyond Accuracy: What Data Quality Means to
Data Consumers.
(2015)}\label{ux10dlanekvir-beyond-accuracy-what-data-quality-means-to-data-consumers.-2015}

\emph{Richard Y. Wang \& Diane M. Strong} Link: {[}5.clanek{]}
(\url{https://www.tandfonline.com/doi/abs/10.1080/07421222.1996.11518099})
-\textgreater{} Kakovost podatkov

\section{Ocena kakovosti podatkov}\label{ocena-kakovosti-podatkov}

\section{Magisterska
naloga-Janez-Bijec--------------------------------}\label{magisterska-naloga-janez-bijec}

Vir/pomoč: Magistrska naloga: Janez Bijec: Statistični kazalniki
kakovosti in učinkoviotsti oskrbe pacientov s koronarno boleznijo

Izpiski: - Kdaj so podatki kakovostni? Kadar so ``fit for purpose'' -
ocenjevanje kakovosti podatkov prilagodit vsaki analizi posebaj glede na
njene zahteve - Avtor v svoji analizi uporabi dva načina: - 1. Na
podlagi modela ocenjevanja in terminologije dimenzij kakovosti
(prilagojen za elektronske zdravstvene podatke) (\emph{skladnost},
\emph{celovitost}, \emph{prepričljivost - verjetnost}) - 2. konceptualni
način

\#\#1.NAČIN--------------------------------------------

\subsection{MODEL OCENJEVANJA IN TERMINOLOGIJE DIMENZIJ
KAKOVOSTI}\label{model-ocenjevanja-in-terminologije-dimenzij-kakovosti}

\begin{enumerate}
\def\labelenumi{\alph{enumi}.}
\tightlist
\item
  SKLADNOST - ali so podatki skladni z definiranimi formati, ustrezajo
  pričakovani obliki in opisom v metapodatkih (poddimenzije: vrednostna
  skladnost, relacijska skladnost, računska skladnost)
\end{enumerate}

\begin{itemize}
\tightlist
\item
  VREDNOSTNA skladnost: spol ima samo dve vrednosti, starost nima
  negativnih vrednosti,\ldots{}
\item
  RELACIJSKA skladnost: omogočanje relacije med tabelami v bazi podatkov
\item
  RAČUNSKA pravilnost izpeljanih spremenljivk
\end{itemize}

\begin{enumerate}
\def\labelenumi{\alph{enumi}.}
\setcounter{enumi}{1}
\item
  CELOVITOST - dimenzija kakovosti, ki ocenjuje prisotnost oz. odsotnost
  (tudi frekvenco podatkov) ne da bi se ozirala na vrednosti, ki jih
  podatki zavzemajo. Ta dimenzija kakovosti nima poddimenzij.
\item
  PREPRIČLJIVOST - VERJETNOST: dimenzija kakovosti, ki osredotoča na
  same vrednosti, ki jih zavzemajo (imajo podatki) in poskuša odgovoriti
  na vprašanje, kako vrednosti podatkov opisujejo realno stanje, koliko
  jim je mogoče zaupati.
\end{enumerate}

Poddimenzije: - UNIKATNA PREPRIČLJIVOST (angl. Uniquness plausibility)
podvajanje podatkov, tam kjer niso smiselni (npr. en identifikator
pacienta je povezan z dvema spoloma ali dvema različnima glavnima
diagnozama za isto hospitalizacijo) - NEČASOVNA PREPRIČLJIVOST (angl.
Atemporal plausibility) Povezana z oceno koliko se vrednosti in
porazdelitve vrednosti skladajo glede na neko splošno vedenje ali glede
na nek zunanju standard. (npr. telesne višine pacientov se porazdeljuje
normalno s povprečjem 170 cim in standardnim odklonom 15cm) - ČASOVNA
PREPLIČLJIVOST (angl. Temporal plausibility) Smiselnost in pravilnost
spreminjanja ali ohranjanja vrednosti spremenljivk glede na čas (npr.
ali je telesna višina pri večini pacientov v adolescenci smiselno
naraščala, odraslega pacienta pa ostala nespremenljena)

\subsection{Ocenjevanje dimenzij
kakovosti-----------------}\label{ocenjevanje-dimenzij-kakovosti}

\subsection{1. VERTIFIKACIJA (angl.
Vertification)}\label{vertifikacija-angl.-vertification}

Temelji na podatkih kot takih na smiselnosti na ujemanju z omejitvami in
strukturo na ekspertnih mnenjih in logičnih povezavah. \#\# 2.
VALIDACIJA (angl. Validation) Temelji na preverjanju podatkov glede na
nek zunanji vir, ki lahko predstavlja zlati standard.

\#\#2.NAČIN-KONCEPTUALNI--------------------------------------------

Avtor poskuša odkriti določene sistemske pomanjklivosti v podatkih.
Ideja Weng in Casey -\textgreater{} predpostavka merljivih konceptov v
medicini

\#OCENA-KAKOVOSTI-PODATKOV-----------------------------

Kakovost podatkov bomo preverili na podlagi šestih najbolj uporabljenih
dimenzij od petnajstih, ki so navedeni v (6.članku). Odločili smo se za
šest dimenzij na podlagi literature in praks.

Glavne dimenzije ocenjevanja kvalitete podatkov, ki jih bomo upoštevali
pri naši analizi so: 1.Natančnost ali točnost\\
2.Celovitost\\
3.Doslednost\\
4.Pravočasnost\\
5.Veljavnost\\
6.Enoličnost

\subsection{Natančnost podatkov}\label{natanux10dnost-podatkov}

``Natančnost podatkov pomeni, kako dobro se informacije ujemajo z
dejstvi iz resničnega sveta. Natančnost podatkov se osredotoča na
pravilnost in zanesljivost podatkov ter z agotavlja, da ti ne vsebujejo
napak, ki bi nastale zaradi nepravilnega vnosa ali pomanjkljive
obdelave. Primarni cilj zagotavljanja točnosti podatkov je prepoznavanje
in odpravljanje napak v podatkovnih vrednostih, kot so napake pri
prepisovanju, podvojeni vnosi in napačne vrednosti. Točnost podatkov se
dosega z uporabo: -- čiščenja podatkov za prepoznavanje in odpravljanje
napak ter neskladnosti v naborih podatkov, kot so odstranjevanje
podvojenih zapisov, popravljanje tipkarskih napak in standardizacija
podatkov; -- preverjanja veljavnosti podatkov s pomočjo vnaprej
določenih pravil ali algoritmov za odkrivanje napak, neskladnosti in
netočnosti v podatkih med postopkom vnosa ali naknadno; -- profiliranja
podatkov za prepoznavanje vzorcev, trendov in nepravilnosti v naborih
podatkov, ki lahko kažejo na netočnosti ali neskladnosti.''

Za naše podatke nimamo zunanjih referenčnih virov s katerim bi lahko
naredili neposredno primerjavo podatkov za preverjanje točnosti
podatkov, ker nimamo neodvisnega vira točnosti podatkov ne moremo
empirično dokazati. Zato bomo v nadaljevanju preverjali logična pravila
in potencialne napačne oz ekstremne vrednosti.

\begin{Shaded}
\begin{Highlighting}[]
\FunctionTok{summary}\NormalTok{(Customer\_Flight\_Activity)}
\end{Highlighting}
\end{Shaded}

\begin{verbatim}
##  Loyalty Number        Year          Month        Total Flights   
##  Min.   :100018   Min.   :2017   Min.   : 1.000   Min.   : 0.000  
##  1st Qu.:327688   1st Qu.:2017   1st Qu.: 4.000   1st Qu.: 0.000  
##  Median :551833   Median :2018   Median : 7.000   Median : 0.000  
##  Mean   :550528   Mean   :2018   Mean   : 6.514   Mean   : 1.295  
##  3rd Qu.:772194   3rd Qu.:2018   3rd Qu.:10.000   3rd Qu.: 2.000  
##  Max.   :999986   Max.   :2018   Max.   :12.000   Max.   :28.000  
##     Distance     Points Accumulated Points Redeemed Dollar Cost Points Redeemed
##  Min.   :    0   Min.   :     0     Min.   :  0.0   Min.   :  0.000            
##  1st Qu.:    0   1st Qu.:     0     1st Qu.:  0.0   1st Qu.:  0.000            
##  Median :    0   Median :     0     Median :  0.0   Median :  0.000            
##  Mean   : 1941   Mean   :  2027     Mean   : 31.3   Mean   :  5.636            
##  3rd Qu.: 3018   3rd Qu.:  3039     3rd Qu.:  0.0   3rd Qu.:  0.000            
##  Max.   :67284   Max.   :100926     Max.   :876.0   Max.   :158.000
\end{verbatim}

\begin{Shaded}
\begin{Highlighting}[]
\FunctionTok{summary}\NormalTok{(Customer\_Loyalty\_History)}
\end{Highlighting}
\end{Shaded}

\begin{verbatim}
##  Loyalty Number     Country            Province             City          
##  Min.   :100018   Length:16737       Length:16737       Length:16737      
##  1st Qu.:326603   Class :character   Class :character   Class :character  
##  Median :550434   Mode  :character   Mode  :character   Mode  :character  
##  Mean   :549736                                                           
##  3rd Qu.:772019                                                           
##  Max.   :999986                                                           
##                                                                           
##  Postal Code           Gender           Education             Salary      
##  Length:16737       Length:16737       Length:16737       Min.   :-58486  
##  Class :character   Class :character   Class :character   1st Qu.: 59246  
##  Mode  :character   Mode  :character   Mode  :character   Median : 73455  
##                                                           Mean   : 79246  
##                                                           3rd Qu.: 88518  
##                                                           Max.   :407228  
##                                                           NA's   :4238    
##  Marital Status     Loyalty Card            CLV        Enrollment Type   
##  Length:16737       Length:16737       Min.   : 1898   Length:16737      
##  Class :character   Class :character   1st Qu.: 3981   Class :character  
##  Mode  :character   Mode  :character   Median : 5780   Mode  :character  
##                                        Mean   : 7989                     
##                                        3rd Qu.: 8941                     
##                                        Max.   :83325                     
##                                                                          
##  Enrollment Year Enrollment Month Cancellation Year Cancellation Month
##  Min.   :2012    Min.   : 1.000   Min.   :2013      Min.   : 1.000    
##  1st Qu.:2014    1st Qu.: 4.000   1st Qu.:2016      1st Qu.: 4.000    
##  Median :2015    Median : 7.000   Median :2017      Median : 7.000    
##  Mean   :2015    Mean   : 6.669   Mean   :2017      Mean   : 6.963    
##  3rd Qu.:2017    3rd Qu.:10.000   3rd Qu.:2018      3rd Qu.:10.000    
##  Max.   :2018    Max.   :12.000   Max.   :2018      Max.   :12.000    
##                                   NA's   :14670     NA's   :14670
\end{verbatim}

Pri spremenljivki `Salary' je videti negativno vrednost, kar ni mogoče,
zaradi obstoja minimalne plače in ker bi negativna vrednost pomenila, da
moremo mi plačevati delodajalcu mesečno plačo.

\begin{Shaded}
\begin{Highlighting}[]
\FunctionTok{ggplot}\NormalTok{(Customer\_Loyalty\_History, }\FunctionTok{aes}\NormalTok{(}\AttributeTok{x =}\NormalTok{ Salary)) }\SpecialCharTok{+}
  \FunctionTok{geom\_histogram}\NormalTok{(}\AttributeTok{binwidth=}\DecValTok{10000}\NormalTok{) }\SpecialCharTok{+}
  \FunctionTok{theme}\NormalTok{(}\AttributeTok{panel.grid =} \FunctionTok{element\_blank}\NormalTok{()) }\SpecialCharTok{+}
  \FunctionTok{labs}\NormalTok{(}
    \AttributeTok{title =} \StringTok{"Porazdelitev spremenljivke Salary"}\NormalTok{,}
    \AttributeTok{x =} \StringTok{"Salary"}\NormalTok{,}
    \AttributeTok{y =} \ConstantTok{NULL}
\NormalTok{  )}
\end{Highlighting}
\end{Shaded}

\begin{verbatim}
## Warning: Removed 4238 rows containing non-finite outside the scale range
## (`stat_bin()`).
\end{verbatim}

\pandocbounded{\includegraphics[keepaspectratio]{Projekt---Airline-Loyalty-Program-Promotion-and-Beyond_files/Projekt---Airline-Loyalty-Program-Promotion-and-Beyond_files/figure-latex/unnamed-chunk-6-1.pdf}}

\subsection{Kakovost podatkov - neurejena Manca
koda}\label{kakovost-podatkov---neurejena-manca-koda}

\begin{Shaded}
\begin{Highlighting}[]
\DocumentationTok{\#\# {-}{-}{-}{-}{-}{-}{-}{-}{-}{-}{-}{-}{-}{-}{-}{-}{-}{-}{-}{-}{-}{-}{-}{-}{-}{-}{-}{-}{-}{-}{-}{-}{-}{-}{-}{-}{-}{-}{-}{-}{-}{-}{-}{-}{-}{-}{-}{-}{-}{-}{-}{-}{-}{-}{-}{-}{-}{-}{-}{-}{-}{-}{-}{-}{-}{-}{-}{-}{-}{-}{-}{-}{-}}
\DocumentationTok{\#\# 2. KAKOVOST PODATKOV}
\DocumentationTok{\#\# {-}{-}{-}{-}{-}{-}{-}{-}{-}{-}{-}{-}{-}{-}{-}{-}{-}{-}{-}{-}{-}{-}{-}{-}{-}{-}{-}{-}{-}{-}{-}{-}{-}{-}{-}{-}{-}{-}{-}{-}{-}{-}{-}{-}{-}{-}{-}{-}{-}{-}{-}{-}{-}{-}{-}{-}{-}{-}{-}{-}{-}{-}{-}{-}{-}{-}{-}{-}{-}{-}{-}{-}{-}}

\NormalTok{flight }\OtherTok{\textless{}{-}}\NormalTok{ Customer\_Flight\_Activity}
\NormalTok{loyalty }\OtherTok{\textless{}{-}}\NormalTok{ Customer\_Loyalty\_History}

\CommentTok{\# 2a. Manjkajoče vrednosti}
\FunctionTok{skim}\NormalTok{(flight)}
\end{Highlighting}
\end{Shaded}

\begin{longtable}[]{@{}ll@{}}
\caption{Data summary}\tabularnewline
\toprule\noalign{}
\endfirsthead
\endhead
\bottomrule\noalign{}
\endlastfoot
Name & flight \\
Number of rows & 392936 \\
Number of columns & 8 \\
\_\_\_\_\_\_\_\_\_\_\_\_\_\_\_\_\_\_\_\_\_\_\_ & \\
Column type frequency: & \\
numeric & 8 \\
\_\_\_\_\_\_\_\_\_\_\_\_\_\_\_\_\_\_\_\_\_\_\_\_ & \\
Group variables & None \\
\end{longtable}

\textbf{Variable type: numeric}

{\def\LTcaptype{none} % do not increment counter
\begin{longtable}[]{@{}
  >{\raggedright\arraybackslash}p{(\linewidth - 20\tabcolsep) * \real{0.2478}}
  >{\raggedleft\arraybackslash}p{(\linewidth - 20\tabcolsep) * \real{0.0885}}
  >{\raggedleft\arraybackslash}p{(\linewidth - 20\tabcolsep) * \real{0.1239}}
  >{\raggedleft\arraybackslash}p{(\linewidth - 20\tabcolsep) * \real{0.0885}}
  >{\raggedleft\arraybackslash}p{(\linewidth - 20\tabcolsep) * \real{0.0885}}
  >{\raggedleft\arraybackslash}p{(\linewidth - 20\tabcolsep) * \real{0.0619}}
  >{\raggedleft\arraybackslash}p{(\linewidth - 20\tabcolsep) * \real{0.0619}}
  >{\raggedleft\arraybackslash}p{(\linewidth - 20\tabcolsep) * \real{0.0619}}
  >{\raggedleft\arraybackslash}p{(\linewidth - 20\tabcolsep) * \real{0.0619}}
  >{\raggedleft\arraybackslash}p{(\linewidth - 20\tabcolsep) * \real{0.0619}}
  >{\raggedright\arraybackslash}p{(\linewidth - 20\tabcolsep) * \real{0.0531}}@{}}
\toprule\noalign{}
\begin{minipage}[b]{\linewidth}\raggedright
skim\_variable
\end{minipage} & \begin{minipage}[b]{\linewidth}\raggedleft
n\_missing
\end{minipage} & \begin{minipage}[b]{\linewidth}\raggedleft
complete\_rate
\end{minipage} & \begin{minipage}[b]{\linewidth}\raggedleft
mean
\end{minipage} & \begin{minipage}[b]{\linewidth}\raggedleft
sd
\end{minipage} & \begin{minipage}[b]{\linewidth}\raggedleft
p0
\end{minipage} & \begin{minipage}[b]{\linewidth}\raggedleft
p25
\end{minipage} & \begin{minipage}[b]{\linewidth}\raggedleft
p50
\end{minipage} & \begin{minipage}[b]{\linewidth}\raggedleft
p75
\end{minipage} & \begin{minipage}[b]{\linewidth}\raggedleft
p100
\end{minipage} & \begin{minipage}[b]{\linewidth}\raggedright
hist
\end{minipage} \\
\midrule\noalign{}
\endhead
\bottomrule\noalign{}
\endlastfoot
Loyalty Number & 0 & 1 & 550527.52 & 258604.58 & 100018 & 327688 &
551833 & 772194 & 999986 & ▇▇▇▇▇ \\
Year & 0 & 1 & 2017.51 & 0.50 & 2017 & 2017 & 2018 & 2018 & 2018 &
▇▁▁▁▇ \\
Month & 0 & 1 & 6.51 & 3.45 & 1 & 4 & 7 & 10 & 12 & ▇▅▅▅▇ \\
Total Flights & 0 & 1 & 1.29 & 1.96 & 0 & 0 & 0 & 2 & 28 & ▇▁▁▁▁ \\
Distance & 0 & 1 & 1941.44 & 3239.98 & 0 & 0 & 0 & 3018 & 67284 &
▇▁▁▁▁ \\
Points Accumulated & 0 & 1 & 2027.17 & 3872.14 & 0 & 0 & 0 & 3039 &
100926 & ▇▁▁▁▁ \\
Points Redeemed & 0 & 1 & 31.30 & 126.65 & 0 & 0 & 0 & 0 & 876 &
▇▁▁▁▁ \\
Dollar Cost Points Redeemed & 0 & 1 & 5.64 & 22.80 & 0 & 0 & 0 & 0 & 158
& ▇▁▁▁▁ \\
\end{longtable}
}

\begin{Shaded}
\begin{Highlighting}[]
\FunctionTok{skim}\NormalTok{(loyalty)}
\end{Highlighting}
\end{Shaded}

\begin{longtable}[]{@{}ll@{}}
\caption{Data summary}\tabularnewline
\toprule\noalign{}
\endfirsthead
\endhead
\bottomrule\noalign{}
\endlastfoot
Name & loyalty \\
Number of rows & 16737 \\
Number of columns & 16 \\
\_\_\_\_\_\_\_\_\_\_\_\_\_\_\_\_\_\_\_\_\_\_\_ & \\
Column type frequency: & \\
character & 9 \\
numeric & 7 \\
\_\_\_\_\_\_\_\_\_\_\_\_\_\_\_\_\_\_\_\_\_\_\_\_ & \\
Group variables & None \\
\end{longtable}

\textbf{Variable type: character}

{\def\LTcaptype{none} % do not increment counter
\begin{longtable}[]{@{}
  >{\raggedright\arraybackslash}p{(\linewidth - 14\tabcolsep) * \real{0.2162}}
  >{\raggedleft\arraybackslash}p{(\linewidth - 14\tabcolsep) * \real{0.1351}}
  >{\raggedleft\arraybackslash}p{(\linewidth - 14\tabcolsep) * \real{0.1892}}
  >{\raggedleft\arraybackslash}p{(\linewidth - 14\tabcolsep) * \real{0.0541}}
  >{\raggedleft\arraybackslash}p{(\linewidth - 14\tabcolsep) * \real{0.0541}}
  >{\raggedleft\arraybackslash}p{(\linewidth - 14\tabcolsep) * \real{0.0811}}
  >{\raggedleft\arraybackslash}p{(\linewidth - 14\tabcolsep) * \real{0.1216}}
  >{\raggedleft\arraybackslash}p{(\linewidth - 14\tabcolsep) * \real{0.1486}}@{}}
\toprule\noalign{}
\begin{minipage}[b]{\linewidth}\raggedright
skim\_variable
\end{minipage} & \begin{minipage}[b]{\linewidth}\raggedleft
n\_missing
\end{minipage} & \begin{minipage}[b]{\linewidth}\raggedleft
complete\_rate
\end{minipage} & \begin{minipage}[b]{\linewidth}\raggedleft
min
\end{minipage} & \begin{minipage}[b]{\linewidth}\raggedleft
max
\end{minipage} & \begin{minipage}[b]{\linewidth}\raggedleft
empty
\end{minipage} & \begin{minipage}[b]{\linewidth}\raggedleft
n\_unique
\end{minipage} & \begin{minipage}[b]{\linewidth}\raggedleft
whitespace
\end{minipage} \\
\midrule\noalign{}
\endhead
\bottomrule\noalign{}
\endlastfoot
Country & 0 & 1 & 6 & 6 & 0 & 1 & 0 \\
Province & 0 & 1 & 5 & 20 & 0 & 11 & 0 \\
City & 0 & 1 & 4 & 14 & 0 & 29 & 0 \\
Postal Code & 0 & 1 & 7 & 7 & 0 & 55 & 0 \\
Gender & 0 & 1 & 4 & 6 & 0 & 2 & 0 \\
Education & 0 & 1 & 6 & 20 & 0 & 5 & 0 \\
Marital Status & 0 & 1 & 6 & 8 & 0 & 3 & 0 \\
Loyalty Card & 0 & 1 & 4 & 6 & 0 & 3 & 0 \\
Enrollment Type & 0 & 1 & 8 & 14 & 0 & 2 & 0 \\
\end{longtable}
}

\textbf{Variable type: numeric}

{\def\LTcaptype{none} % do not increment counter
\begin{longtable}[]{@{}
  >{\raggedright\arraybackslash}p{(\linewidth - 20\tabcolsep) * \real{0.1597}}
  >{\raggedleft\arraybackslash}p{(\linewidth - 20\tabcolsep) * \real{0.0840}}
  >{\raggedleft\arraybackslash}p{(\linewidth - 20\tabcolsep) * \real{0.1176}}
  >{\raggedleft\arraybackslash}p{(\linewidth - 20\tabcolsep) * \real{0.0840}}
  >{\raggedleft\arraybackslash}p{(\linewidth - 20\tabcolsep) * \real{0.0840}}
  >{\raggedleft\arraybackslash}p{(\linewidth - 20\tabcolsep) * \real{0.0840}}
  >{\raggedleft\arraybackslash}p{(\linewidth - 20\tabcolsep) * \real{0.0840}}
  >{\raggedleft\arraybackslash}p{(\linewidth - 20\tabcolsep) * \real{0.0840}}
  >{\raggedleft\arraybackslash}p{(\linewidth - 20\tabcolsep) * \real{0.0840}}
  >{\raggedleft\arraybackslash}p{(\linewidth - 20\tabcolsep) * \real{0.0840}}
  >{\raggedright\arraybackslash}p{(\linewidth - 20\tabcolsep) * \real{0.0504}}@{}}
\toprule\noalign{}
\begin{minipage}[b]{\linewidth}\raggedright
skim\_variable
\end{minipage} & \begin{minipage}[b]{\linewidth}\raggedleft
n\_missing
\end{minipage} & \begin{minipage}[b]{\linewidth}\raggedleft
complete\_rate
\end{minipage} & \begin{minipage}[b]{\linewidth}\raggedleft
mean
\end{minipage} & \begin{minipage}[b]{\linewidth}\raggedleft
sd
\end{minipage} & \begin{minipage}[b]{\linewidth}\raggedleft
p0
\end{minipage} & \begin{minipage}[b]{\linewidth}\raggedleft
p25
\end{minipage} & \begin{minipage}[b]{\linewidth}\raggedleft
p50
\end{minipage} & \begin{minipage}[b]{\linewidth}\raggedleft
p75
\end{minipage} & \begin{minipage}[b]{\linewidth}\raggedleft
p100
\end{minipage} & \begin{minipage}[b]{\linewidth}\raggedright
hist
\end{minipage} \\
\midrule\noalign{}
\endhead
\bottomrule\noalign{}
\endlastfoot
Loyalty Number & 0 & 1.00 & 549735.88 & 258912.13 & 100018.00 &
326603.00 & 550434.00 & 772019.00 & 999986.00 & ▇▇▇▇▇ \\
Salary & 4238 & 0.75 & 79245.61 & 35008.30 & -58486.00 & 59246.50 &
73455.00 & 88517.50 & 407228.00 & ▁▇▁▁▁ \\
CLV & 0 & 1.00 & 7988.90 & 6860.98 & 1898.01 & 3980.84 & 5780.18 &
8940.58 & 83325.38 & ▇▁▁▁▁ \\
Enrollment Year & 0 & 1.00 & 2015.25 & 1.98 & 2012.00 & 2014.00 &
2015.00 & 2017.00 & 2018.00 & ▆▃▃▃▇ \\
Enrollment Month & 0 & 1.00 & 6.67 & 3.40 & 1.00 & 4.00 & 7.00 & 10.00 &
12.00 & ▇▅▅▅▇ \\
Cancellation Year & 14670 & 0.12 & 2016.50 & 1.38 & 2013.00 & 2016.00 &
2017.00 & 2018.00 & 2018.00 & ▃▃▅▆▇ \\
Cancellation Month & 14670 & 0.12 & 6.96 & 3.46 & 1.00 & 4.00 & 7.00 &
10.00 & 12.00 & ▆▃▅▅▇ \\
\end{longtable}
}

\begin{Shaded}
\begin{Highlighting}[]
\CommentTok{\# Delež NA po stolpcu}
\FunctionTok{map\_dbl}\NormalTok{(loyalty, }\SpecialCharTok{\textasciitilde{}}\FunctionTok{mean}\NormalTok{(}\FunctionTok{is.na}\NormalTok{(.))) }\SpecialCharTok{\%\textgreater{}\%} \FunctionTok{sort}\NormalTok{(}\AttributeTok{decreasing =} \ConstantTok{TRUE}\NormalTok{)}
\end{Highlighting}
\end{Shaded}

\begin{verbatim}
##  Cancellation Year Cancellation Month             Salary     Loyalty Number 
##          0.8765012          0.8765012          0.2532114          0.0000000 
##            Country           Province               City        Postal Code 
##          0.0000000          0.0000000          0.0000000          0.0000000 
##             Gender          Education     Marital Status       Loyalty Card 
##          0.0000000          0.0000000          0.0000000          0.0000000 
##                CLV    Enrollment Type    Enrollment Year   Enrollment Month 
##          0.0000000          0.0000000          0.0000000          0.0000000
\end{verbatim}

V izpisanih tabelah lahko vidimo, da podatki o letih nimajo manjkajočih
vrednosti ter hitre histogramske porazdelitve posamičnih spremenljivk.
Medtem ko pri podatkih o članih programa opazimo, da imamo manjkajoče
vrednosti pri spremenljivki Salary in pa seveda pri spremenljivkama
Cancellation Year in Month, ki pa niso manjkajoči podatki, vendar so to
še aktualni člani in podatka enostavno ni. Tudi v teh podatkih lahko
pogledamo histograme porazdelitev številskih spremenljivk.

V nadaljnem smo pogledali le še delež manjkajočih vrednosti in opazimo,
da je 87.6\% vseh članov še vedno aktivnih, kar je super odstotek in da
nam pri spremenljivki o plačah manjka 25\% vrednosti. To so dejankse
manjkajoče vrednosti, ki si jih želimo še malo bolj natančno pogledati,
ali obstaja kakšen vzorec manjkajočnosti.

\begin{Shaded}
\begin{Highlighting}[]
\CommentTok{\# 2b. Dvojniki po ključu}
\NormalTok{loyalty }\SpecialCharTok{\%\textgreater{}\%} \FunctionTok{count}\NormalTok{(}\StringTok{\textasciigrave{}}\AttributeTok{Loyalty Number}\StringTok{\textasciigrave{}}\NormalTok{) }\SpecialCharTok{\%\textgreater{}\%} \FunctionTok{filter}\NormalTok{(n }\SpecialCharTok{\textgreater{}} \DecValTok{1}\NormalTok{) }\CommentTok{\# naj bo prazno {-} 1 vrstica na stranko}
\end{Highlighting}
\end{Shaded}

\begin{verbatim}
## # A tibble: 0 x 2
## # i 2 variables: Loyalty Number <dbl>, n <int>
\end{verbatim}

\begin{Shaded}
\begin{Highlighting}[]
\NormalTok{flight  }\SpecialCharTok{\%\textgreater{}\%} \FunctionTok{count}\NormalTok{(}\StringTok{\textasciigrave{}}\AttributeTok{Loyalty Number}\StringTok{\textasciigrave{}}\NormalTok{, Year, Month) }\SpecialCharTok{\%\textgreater{}\%} \FunctionTok{filter}\NormalTok{(n }\SpecialCharTok{\textgreater{}} \DecValTok{1}\NormalTok{) }\CommentTok{\# naj bo prazno}
\end{Highlighting}
\end{Shaded}

\begin{verbatim}
## # A tibble: 3,847 x 4
##    `Loyalty Number`  Year Month     n
##               <dbl> <dbl> <dbl> <int>
##  1           101902  2017     1     2
##  2           101902  2017     2     2
##  3           101902  2017     3     2
##  4           101902  2017     4     2
##  5           101902  2017     5     2
##  6           101902  2017     6     2
##  7           101902  2017     7     2
##  8           101902  2017     8     2
##  9           101902  2017     9     2
## 10           101902  2017    10     2
## # i 3,837 more rows
\end{verbatim}

Velik kazatelj kvalitetnih podatkov je tudi unikatnost vnosov, ker imamo
namreč osebne podatke v eni datoteki, se tam dvojniki po številki
zvestobe ne bi smeli prikazovati, prav tako bi v datoteki z leti vsak
član moral imeti za vsak mesec v letu le en vnos.

V zgornjem izpisu vidimo, da dvojnih vnosov o strankah ni, je pa kar
veliko dvojnih vnosov v tabeli z leti, kar je lahko slab kazatelj
kakovosti podatkov, ki pa zahteva nadaljni vpogled.

\begin{Shaded}
\begin{Highlighting}[]
\CommentTok{\# 2c. Ali imajo vsi poleti veljavno stranko v loyalty tabeli? (referenčna integriteta)}
\FunctionTok{anti\_join}\NormalTok{(flight, loyalty, }\AttributeTok{by =} \StringTok{"Loyalty Number"}\NormalTok{) }\SpecialCharTok{\%\textgreater{}\%} \FunctionTok{nrow}\NormalTok{()}
\end{Highlighting}
\end{Shaded}

\begin{verbatim}
## [1] 0
\end{verbatim}

\begin{Shaded}
\begin{Highlighting}[]
\CommentTok{\# 2d. Osamelci / nesmiselne vrednosti}
\FunctionTok{summary}\NormalTok{(flight}\SpecialCharTok{$}\StringTok{\textasciigrave{}}\AttributeTok{Total Flights}\StringTok{\textasciigrave{}}\NormalTok{)}
\end{Highlighting}
\end{Shaded}

\begin{verbatim}
##    Min. 1st Qu.  Median    Mean 3rd Qu.    Max. 
##   0.000   0.000   0.000   1.295   2.000  28.000
\end{verbatim}

\begin{Shaded}
\begin{Highlighting}[]
\FunctionTok{summary}\NormalTok{(flight}\SpecialCharTok{$}\NormalTok{Distance)}
\end{Highlighting}
\end{Shaded}

\begin{verbatim}
##    Min. 1st Qu.  Median    Mean 3rd Qu.    Max. 
##       0       0       0    1941    3018   67284
\end{verbatim}

\begin{Shaded}
\begin{Highlighting}[]
\NormalTok{flight }\SpecialCharTok{\%\textgreater{}\%} \FunctionTok{filter}\NormalTok{(}\StringTok{\textasciigrave{}}\AttributeTok{Total Flights}\StringTok{\textasciigrave{}} \SpecialCharTok{==} \DecValTok{0} \SpecialCharTok{\&}\NormalTok{ Distance }\SpecialCharTok{\textgreater{}} \DecValTok{0}\NormalTok{) }\CommentTok{\# nekonsistentnost}
\end{Highlighting}
\end{Shaded}

\begin{verbatim}
## # A tibble: 0 x 8
## # i 8 variables: Loyalty Number <dbl>, Year <dbl>, Month <dbl>,
## #   Total Flights <dbl>, Distance <dbl>, Points Accumulated <dbl>,
## #   Points Redeemed <dbl>, Dollar Cost Points Redeemed <dbl>
\end{verbatim}

\begin{Shaded}
\begin{Highlighting}[]
\NormalTok{flight }\SpecialCharTok{\%\textgreater{}\%} \FunctionTok{filter}\NormalTok{(}\StringTok{\textasciigrave{}}\AttributeTok{Total Flights}\StringTok{\textasciigrave{}} \SpecialCharTok{\textgreater{}} \DecValTok{0} \SpecialCharTok{\&}\NormalTok{ Distance }\SpecialCharTok{==} \DecValTok{0}\NormalTok{)}
\end{Highlighting}
\end{Shaded}

\begin{verbatim}
## # A tibble: 0 x 8
## # i 8 variables: Loyalty Number <dbl>, Year <dbl>, Month <dbl>,
## #   Total Flights <dbl>, Distance <dbl>, Points Accumulated <dbl>,
## #   Points Redeemed <dbl>, Dollar Cost Points Redeemed <dbl>
\end{verbatim}

\begin{Shaded}
\begin{Highlighting}[]
\FunctionTok{summary}\NormalTok{(loyalty}\SpecialCharTok{$}\NormalTok{Salary)}
\end{Highlighting}
\end{Shaded}

\begin{verbatim}
##    Min. 1st Qu.  Median    Mean 3rd Qu.    Max.    NA's 
##  -58486   59246   73455   79246   88518  407228    4238
\end{verbatim}

\begin{Shaded}
\begin{Highlighting}[]
\NormalTok{loyalty }\SpecialCharTok{\%\textgreater{}\%} \FunctionTok{filter}\NormalTok{(Salary }\SpecialCharTok{\textless{}} \DecValTok{0}\NormalTok{) }\CommentTok{\# negativna plača = napaka?}
\end{Highlighting}
\end{Shaded}

\begin{verbatim}
## # A tibble: 20 x 16
##    `Loyalty Number` Country Province City  `Postal Code` Gender Education Salary
##               <dbl> <chr>   <chr>    <chr> <chr>         <chr>  <chr>      <dbl>
##  1           542976 Canada  Quebec   Mont~ H2Y 4R4       Male   High Sch~ -49830
##  2           959977 Canada  British~ Vanc~ V5R 1W3       Female Bachelor  -12497
##  3           232755 Canada  British~ Vanc~ V1E 4R6       Female Bachelor  -46683
##  4           525245 Canada  British~ Vict~ V10 6T5       Male   Bachelor  -45962
##  5           603070 Canada  British~ West~ V6V 8Z3       Female Bachelor  -19325
##  6           491242 Canada  British~ Daws~ U5I 4F1       Male   Bachelor  -43234
##  7           115505 Canada  Newfoun~ St. ~ A1C 6H9       Male   Bachelor  -10605
##  8           430398 Canada  Newfoun~ St. ~ A1C 6H9       Male   Bachelor  -17534
##  9           152016 Canada  Ontario  Toro~ P1J 8T7       Female Bachelor  -58486
## 10           194065 Canada  Ontario  Sudb~ M5V 1G5       Female Bachelor  -31911
## 11           212128 Canada  Ontario  Toro~ P2T 6G3       Male   Bachelor  -49001
## 12           790475 Canada  Ontario  Tren~ K8V 4B2       Female Bachelor  -34079
## 13           366599 Canada  Ontario  Toro~ M1R 4K3       Female Bachelor   -9081
## 14           436087 Canada  Quebec   Mont~ H2T 9K8       Male   Bachelor  -46470
## 15           364596 Canada  Quebec   Trem~ H5Y 2S9       Female Bachelor  -26322
## 16           239955 Canada  Quebec   Hull  J8Y 3Z5       Female Bachelor  -47310
## 17           347013 Canada  Quebec   Queb~ G1B 3L5       Male   Bachelor  -39503
## 18           729561 Canada  Quebec   Queb~ G1B 3L5       Female Bachelor  -19332
## 19           734647 Canada  Saskatc~ Regi~ S1J 3C5       Male   Bachelor  -46303
## 20           906428 Canada  Yukon    Whit~ Y2K 6R0       Male   Bachelor  -57297
## # i 8 more variables: `Marital Status` <chr>, `Loyalty Card` <chr>, CLV <dbl>,
## #   `Enrollment Type` <chr>, `Enrollment Year` <dbl>, `Enrollment Month` <dbl>,
## #   `Cancellation Year` <dbl>, `Cancellation Month` <dbl>
\end{verbatim}

\begin{Shaded}
\begin{Highlighting}[]
\CommentTok{\# 2e. Konsistentnost kategorij}
\NormalTok{loyalty }\SpecialCharTok{\%\textgreater{}\%} \FunctionTok{count}\NormalTok{(}\StringTok{\textasciigrave{}}\AttributeTok{Enrollment Type}\StringTok{\textasciigrave{}}\NormalTok{)}
\end{Highlighting}
\end{Shaded}

\begin{verbatim}
## # A tibble: 2 x 2
##   `Enrollment Type`     n
##   <chr>             <int>
## 1 2018 Promotion      971
## 2 Standard          15766
\end{verbatim}

\begin{Shaded}
\begin{Highlighting}[]
\NormalTok{loyalty }\SpecialCharTok{\%\textgreater{}\%} \FunctionTok{count}\NormalTok{(}\StringTok{\textasciigrave{}}\AttributeTok{Loyalty Card}\StringTok{\textasciigrave{}}\NormalTok{)}
\end{Highlighting}
\end{Shaded}

\begin{verbatim}
## # A tibble: 3 x 2
##   `Loyalty Card`     n
##   <chr>          <int>
## 1 Aurora          3429
## 2 Nova            5671
## 3 Star            7637
\end{verbatim}

\begin{Shaded}
\begin{Highlighting}[]
\NormalTok{loyalty }\SpecialCharTok{\%\textgreater{}\%} \FunctionTok{count}\NormalTok{(Education)}
\end{Highlighting}
\end{Shaded}

\begin{verbatim}
## # A tibble: 5 x 2
##   Education                n
##   <chr>                <int>
## 1 Bachelor             10475
## 2 College               4238
## 3 Doctor                 734
## 4 High School or Below   782
## 5 Master                 508
\end{verbatim}

\begin{Shaded}
\begin{Highlighting}[]
\NormalTok{loyalty }\SpecialCharTok{\%\textgreater{}\%} \FunctionTok{count}\NormalTok{(}\StringTok{\textasciigrave{}}\AttributeTok{Marital Status}\StringTok{\textasciigrave{}}\NormalTok{)}
\end{Highlighting}
\end{Shaded}

\begin{verbatim}
## # A tibble: 3 x 2
##   `Marital Status`     n
##   <chr>            <int>
## 1 Divorced          2518
## 2 Married           9735
## 3 Single            4484
\end{verbatim}

\begin{Shaded}
\begin{Highlighting}[]
\NormalTok{loyalty }\SpecialCharTok{\%\textgreater{}\%} \FunctionTok{count}\NormalTok{(Gender)}
\end{Highlighting}
\end{Shaded}

\begin{verbatim}
## # A tibble: 2 x 2
##   Gender     n
##   <chr>  <int>
## 1 Female  8410
## 2 Male    8327
\end{verbatim}

Pogledali smo še, da imajo vsi vnosi v letih veljavno stranko, da so
podatki o poletih v enem mesecu in razdalji smiselni, ter da ne obstaja
vnos, ki bi imel nekonsistentno povezano razdaljo in število letov. Prav
tako smo pogledali, da so plače članov programa smiselne, kjer smo
opazili nekaj negativnih plač, ki jih bomo izbrisali iz vnosa in
obravnavali kot manjkajoče podatke, saj ne moremo vedeti, ali so samo
predznaki napaka ali je cel vnos napačen.

\subsection{Popravek podatkov - neurejena Manca
koda}\label{popravek-podatkov---neurejena-manca-koda}

Tole mi ni čist všeč še. Cilji so: - odstraniti vse lete nepopolnih
dvojnikov, ker ne vemo, kateri je pravi - odstraniti popolne dvojnike,
ker lahko zaupamo, da so vseeno pravilne vrednosti - odstrani negativne
plače

\begin{Shaded}
\begin{Highlighting}[]
\DocumentationTok{\#\# {-}{-}{-}{-}{-}{-}{-}{-}{-}{-}{-}{-}{-}{-}{-}{-}{-}{-}{-}{-}{-}{-}{-}{-}{-}{-}{-}{-}{-}{-}{-}{-}{-}{-}{-}{-}{-}{-}{-}{-}{-}{-}{-}{-}{-}{-}{-}{-}{-}{-}{-}{-}{-}{-}{-}{-}{-}{-}{-}{-}{-}{-}{-}{-}{-}{-}{-}{-}{-}{-}{-}{-}{-}}
\DocumentationTok{\#\# 3. ČIŠČENJE PODATKOV}
\DocumentationTok{\#\# {-}{-}{-}{-}{-}{-}{-}{-}{-}{-}{-}{-}{-}{-}{-}{-}{-}{-}{-}{-}{-}{-}{-}{-}{-}{-}{-}{-}{-}{-}{-}{-}{-}{-}{-}{-}{-}{-}{-}{-}{-}{-}{-}{-}{-}{-}{-}{-}{-}{-}{-}{-}{-}{-}{-}{-}{-}{-}{-}{-}{-}{-}{-}{-}{-}{-}{-}{-}{-}{-}{-}{-}{-}}

\DocumentationTok{\#\# 3.0 Čiščenje podvojenih vrstic v flight\_activity}
\DocumentationTok{\#\# {-}{-}{-}{-}{-}{-}{-}{-}{-}{-}{-}{-}{-}{-}{-}{-}{-}{-}{-}{-}{-}{-}{-}{-}{-}{-}{-}{-}{-}{-}{-}{-}{-}{-}{-}{-}{-}{-}{-}{-}{-}{-}{-}{-}{-}{-}{-}{-}{-}{-}{-}{-}{-}{-}{-}{-}{-}{-}{-}{-}{-}{-}{-}{-}{-}{-}{-}{-}{-}{-}{-}{-}{-}}
\DocumentationTok{\#\# Pravilo:}
\DocumentationTok{\#\#   a) Če imajo vrstice z enakim ključem (Loyalty Number, Year, Month) RAZLIČNE vrednosti drugje {-}\textgreater{} ne vemo, katera je prava, zato odstranimo VSE.}
\DocumentationTok{\#\#   b) Če so vrstice z enakim ključem POPOLNI dvojniki (identične povsod) {-}\textgreater{} zaupamo vrednosti in obdržimo eno kopijo.}

\NormalTok{key\_cols }\OtherTok{\textless{}{-}} \FunctionTok{c}\NormalTok{(}\StringTok{"Loyalty Number"}\NormalTok{, }\StringTok{"Year"}\NormalTok{, }\StringTok{"Month"}\NormalTok{)}

\CommentTok{\# Za vsako skupino preveri velikost skupine in št. dejansko različnih vrstic znotraj nje}
\NormalTok{key\_dup\_summary }\OtherTok{\textless{}{-}}\NormalTok{ flight }\SpecialCharTok{\%\textgreater{}\%}
  \FunctionTok{group\_by}\NormalTok{(}\FunctionTok{across}\NormalTok{(}\FunctionTok{all\_of}\NormalTok{(key\_cols))) }\SpecialCharTok{\%\textgreater{}\%}
  \FunctionTok{summarise}\NormalTok{(}
    \AttributeTok{group\_size        =} \FunctionTok{n}\NormalTok{(),}
    \AttributeTok{n\_distinct\_rows    =} \FunctionTok{n\_distinct}\NormalTok{(}\FunctionTok{pick}\NormalTok{(}\FunctionTok{everything}\NormalTok{())),}
    \AttributeTok{.groups =} \StringTok{"drop"}
\NormalTok{  ) }\SpecialCharTok{\%\textgreater{}\%}
  \FunctionTok{filter}\NormalTok{(group\_size }\SpecialCharTok{\textgreater{}} \DecValTok{1}\NormalTok{)}

\CommentTok{\# Ključi, kjer je skupina POPOLN dvojnik (n\_distinct\_rows == 1 kljub group\_size \textgreater{} 1)}
\NormalTok{full\_dup\_keys }\OtherTok{\textless{}{-}}\NormalTok{ key\_dup\_summary }\SpecialCharTok{\%\textgreater{}\%}
  \FunctionTok{filter}\NormalTok{(n\_distinct\_rows }\SpecialCharTok{==} \DecValTok{1}\NormalTok{) }\SpecialCharTok{\%\textgreater{}\%}
  \FunctionTok{select}\NormalTok{(}\FunctionTok{all\_of}\NormalTok{(key\_cols))}

\CommentTok{\# Ključi, kjer skupina NI popoln dvojnik (vrstice se med sabo razlikujejo)}
\NormalTok{partial\_dup\_keys }\OtherTok{\textless{}{-}}\NormalTok{ key\_dup\_summary }\SpecialCharTok{\%\textgreater{}\%}
  \FunctionTok{filter}\NormalTok{(n\_distinct\_rows }\SpecialCharTok{\textgreater{}} \DecValTok{1}\NormalTok{) }\SpecialCharTok{\%\textgreater{}\%}
  \FunctionTok{select}\NormalTok{(}\FunctionTok{all\_of}\NormalTok{(key\_cols))}

\FunctionTok{nrow}\NormalTok{(full\_dup\_keys)      }\CommentTok{\# št. ključev s popolnimi dvojniki (obdržimo 1 kopijo)}
\end{Highlighting}
\end{Shaded}

\begin{verbatim}
## [1] 1904
\end{verbatim}

\begin{Shaded}
\begin{Highlighting}[]
\FunctionTok{nrow}\NormalTok{(partial\_dup\_keys)   }\CommentTok{\# št. ključev z nepopolnimi/sumljivimi dvojniki (odstranimo VSE)}
\end{Highlighting}
\end{Shaded}

\begin{verbatim}
## [1] 1943
\end{verbatim}

\begin{Shaded}
\begin{Highlighting}[]
\CommentTok{\# Vpogled v primer nepopolnih dvojnikov}
\NormalTok{flight }\SpecialCharTok{\%\textgreater{}\%}
  \FunctionTok{semi\_join}\NormalTok{(partial\_dup\_keys, }\AttributeTok{by =}\NormalTok{ key\_cols) }\SpecialCharTok{\%\textgreater{}\%}
  \FunctionTok{arrange}\NormalTok{(}\StringTok{\textasciigrave{}}\AttributeTok{Loyalty Number}\StringTok{\textasciigrave{}}\NormalTok{, Year, Month) }\SpecialCharTok{\%\textgreater{}\%}
  \FunctionTok{head}\NormalTok{(}\DecValTok{10}\NormalTok{)}
\end{Highlighting}
\end{Shaded}

\begin{verbatim}
## # A tibble: 10 x 8
##    `Loyalty Number`  Year Month `Total Flights` Distance `Points Accumulated`
##               <dbl> <dbl> <dbl>           <dbl>    <dbl>                <dbl>
##  1           101902  2017     4               1      667                  667
##  2           101902  2017     4               1     2499                 2499
##  3           101902  2017     5               3     7263                 7263
##  4           101902  2017     5               2     3762                 3762
##  5           101902  2017     6               3     4971                 4971
##  6           101902  2017     6               3     5763                 5763
##  7           101902  2017     8               5     3825                 3825
##  8           101902  2017     8               4     8952                 8952
##  9           101902  2017     9               2     3734                 3734
## 10           101902  2017     9               1     2187                 2187
## # i 2 more variables: `Points Redeemed` <dbl>,
## #   `Dollar Cost Points Redeemed` <dbl>
\end{verbatim}

\begin{Shaded}
\begin{Highlighting}[]
\CommentTok{\# Odstranitev:}
\CommentTok{\#   (a) vse vrstice, ki pripadajo nepopolnim/sumljivim dvojnikom}
\CommentTok{\#   (b) med preostalimi vrsticami odstranimo odvečne kopije popolnih dvojnikov}
\NormalTok{flight\_clean }\OtherTok{\textless{}{-}}\NormalTok{ flight }\SpecialCharTok{\%\textgreater{}\%}
  \FunctionTok{anti\_join}\NormalTok{(partial\_dup\_keys, }\AttributeTok{by =}\NormalTok{ key\_cols) }\SpecialCharTok{\%\textgreater{}\%}
  \FunctionTok{distinct}\NormalTok{()}

\CommentTok{\# Poročilo o čiščenju}
\NormalTok{n\_total            }\OtherTok{\textless{}{-}} \FunctionTok{nrow}\NormalTok{(flight)}
\NormalTok{n\_removed\_partial  }\OtherTok{\textless{}{-}}\NormalTok{ flight }\SpecialCharTok{\%\textgreater{}\%} \FunctionTok{semi\_join}\NormalTok{(partial\_dup\_keys, }\AttributeTok{by =}\NormalTok{ key\_cols) }\SpecialCharTok{\%\textgreater{}\%} \FunctionTok{nrow}\NormalTok{()}
\NormalTok{n\_removed\_full\_dup }\OtherTok{\textless{}{-}}\NormalTok{ n\_total }\SpecialCharTok{{-}}\NormalTok{ n\_removed\_partial }\SpecialCharTok{{-}} \FunctionTok{nrow}\NormalTok{(flight\_clean)}

\FunctionTok{cat}\NormalTok{(}\StringTok{"Vrstic prej:"}\NormalTok{, n\_total, }\StringTok{"}\SpecialCharTok{\textbackslash{}n}\StringTok{"}\NormalTok{)}
\end{Highlighting}
\end{Shaded}

\begin{verbatim}
## Vrstic prej: 392936
\end{verbatim}

\begin{Shaded}
\begin{Highlighting}[]
\FunctionTok{cat}\NormalTok{(}\StringTok{"Odstranjenih (nepopolni dvojniki {-} vse vrstice skupine):"}\NormalTok{, n\_removed\_partial, }\StringTok{"}\SpecialCharTok{\textbackslash{}n}\StringTok{"}\NormalTok{)}
\end{Highlighting}
\end{Shaded}

\begin{verbatim}
## Odstranjenih (nepopolni dvojniki - vse vrstice skupine): 3894
\end{verbatim}

\begin{Shaded}
\begin{Highlighting}[]
\FunctionTok{cat}\NormalTok{(}\StringTok{"Odstranjenih (odvečne kopije popolnih dvojnikov):"}\NormalTok{, n\_removed\_full\_dup, }\StringTok{"}\SpecialCharTok{\textbackslash{}n}\StringTok{"}\NormalTok{)}
\end{Highlighting}
\end{Shaded}

\begin{verbatim}
## Odstranjenih (odvečne kopije popolnih dvojnikov): 1920
\end{verbatim}

\begin{Shaded}
\begin{Highlighting}[]
\FunctionTok{cat}\NormalTok{(}\StringTok{"Vrstic po čiščenju:"}\NormalTok{, }\FunctionTok{nrow}\NormalTok{(flight\_clean), }\StringTok{"}\SpecialCharTok{\textbackslash{}n}\StringTok{"}\NormalTok{)}
\end{Highlighting}
\end{Shaded}

\begin{verbatim}
## Vrstic po čiščenju: 387122
\end{verbatim}

\begin{Shaded}
\begin{Highlighting}[]
\CommentTok{\# Kontrola {-} noben ključ naj se več ne podvaja}
\NormalTok{flight\_clean }\SpecialCharTok{\%\textgreater{}\%} \FunctionTok{count}\NormalTok{(}\FunctionTok{across}\NormalTok{(}\FunctionTok{all\_of}\NormalTok{(key\_cols))) }\SpecialCharTok{\%\textgreater{}\%} \FunctionTok{filter}\NormalTok{(n }\SpecialCharTok{\textgreater{}} \DecValTok{1}\NormalTok{)}
\end{Highlighting}
\end{Shaded}

\begin{verbatim}
## # A tibble: 0 x 4
## # i 4 variables: Loyalty Number <dbl>, Year <dbl>, Month <dbl>, n <int>
\end{verbatim}

\begin{Shaded}
\begin{Highlighting}[]
\DocumentationTok{\#\# 3.1 Odstranitev negativnih plač v loyalty\_history}
\DocumentationTok{\#\# {-}{-}{-}{-}{-}{-}{-}{-}{-}{-}{-}{-}{-}{-}{-}{-}{-}{-}{-}{-}{-}{-}{-}{-}{-}{-}{-}{-}{-}{-}{-}{-}{-}{-}{-}{-}{-}{-}{-}{-}{-}{-}{-}{-}{-}{-}{-}{-}{-}{-}{-}{-}{-}{-}{-}{-}{-}{-}{-}{-}{-}{-}{-}{-}{-}{-}{-}{-}{-}{-}{-}{-}{-}}

\NormalTok{n\_negative\_salary }\OtherTok{\textless{}{-}}\NormalTok{ loyalty }\SpecialCharTok{\%\textgreater{}\%} \FunctionTok{filter}\NormalTok{(Salary }\SpecialCharTok{\textless{}} \DecValTok{0}\NormalTok{) }\SpecialCharTok{\%\textgreater{}\%} \FunctionTok{nrow}\NormalTok{()}
\NormalTok{n\_negative\_salary   }\CommentTok{\# 20}
\end{Highlighting}
\end{Shaded}

\begin{verbatim}
## [1] 20
\end{verbatim}

\begin{Shaded}
\begin{Highlighting}[]
\NormalTok{loyalty\_clean }\OtherTok{\textless{}{-}}\NormalTok{ loyalty }\SpecialCharTok{\%\textgreater{}\%}
  \FunctionTok{mutate}\NormalTok{(}\AttributeTok{Salary =} \FunctionTok{if\_else}\NormalTok{(Salary }\SpecialCharTok{\textless{}} \DecValTok{0}\NormalTok{, }\ConstantTok{NA\_real\_}\NormalTok{, Salary))}

\DocumentationTok{\#\# 3.2 Agregacija po stranki (na podlagi flight\_clean)}
\DocumentationTok{\#\# {-}{-}{-}{-}{-}{-}{-}{-}{-}{-}{-}{-}{-}{-}{-}{-}{-}{-}{-}{-}{-}{-}{-}{-}{-}{-}{-}{-}{-}{-}{-}{-}{-}{-}{-}{-}{-}{-}{-}{-}{-}{-}{-}{-}{-}{-}{-}{-}{-}{-}{-}{-}{-}{-}{-}{-}{-}{-}{-}{-}{-}{-}{-}{-}{-}{-}{-}{-}{-}{-}{-}{-}{-}}

\NormalTok{customer\_activity }\OtherTok{\textless{}{-}}\NormalTok{ flight\_clean }\SpecialCharTok{\%\textgreater{}\%}
  \FunctionTok{group\_by}\NormalTok{(}\StringTok{\textasciigrave{}}\AttributeTok{Loyalty Number}\StringTok{\textasciigrave{}}\NormalTok{) }\SpecialCharTok{\%\textgreater{}\%}
  \FunctionTok{summarise}\NormalTok{(}
    \AttributeTok{total\_flights\_sum    =} \FunctionTok{sum}\NormalTok{(}\StringTok{\textasciigrave{}}\AttributeTok{Total Flights}\StringTok{\textasciigrave{}}\NormalTok{, }\AttributeTok{na.rm =} \ConstantTok{TRUE}\NormalTok{),}
    \AttributeTok{total\_distance       =} \FunctionTok{sum}\NormalTok{(Distance, }\AttributeTok{na.rm =} \ConstantTok{TRUE}\NormalTok{),}
    \AttributeTok{points\_accumulated   =} \FunctionTok{sum}\NormalTok{(}\StringTok{\textasciigrave{}}\AttributeTok{Points Accumulated}\StringTok{\textasciigrave{}}\NormalTok{, }\AttributeTok{na.rm =} \ConstantTok{TRUE}\NormalTok{),}
    \AttributeTok{points\_redeemed      =} \FunctionTok{sum}\NormalTok{(}\StringTok{\textasciigrave{}}\AttributeTok{Points Redeemed}\StringTok{\textasciigrave{}}\NormalTok{, }\AttributeTok{na.rm =} \ConstantTok{TRUE}\NormalTok{),}
    \AttributeTok{dollar\_cost\_redeemed =} \FunctionTok{sum}\NormalTok{(}\StringTok{\textasciigrave{}}\AttributeTok{Dollar Cost Points Redeemed}\StringTok{\textasciigrave{}}\NormalTok{, }\AttributeTok{na.rm =} \ConstantTok{TRUE}\NormalTok{),}
    \AttributeTok{active\_months        =} \FunctionTok{n\_distinct}\NormalTok{(}\FunctionTok{paste}\NormalTok{(Year, Month)),}
    \AttributeTok{.groups =} \StringTok{"drop"}
\NormalTok{  ) }\SpecialCharTok{\%\textgreater{}\%}
  \FunctionTok{mutate}\NormalTok{(}
    \AttributeTok{redemption\_rate       =} \FunctionTok{if\_else}\NormalTok{(points\_accumulated }\SpecialCharTok{\textgreater{}} \DecValTok{0}\NormalTok{,}
\NormalTok{                                     points\_redeemed }\SpecialCharTok{/}\NormalTok{ points\_accumulated, }\DecValTok{0}\NormalTok{),}
    \AttributeTok{avg\_flights\_per\_month =}\NormalTok{ total\_flights\_sum }\SpecialCharTok{/}\NormalTok{ active\_months}
\NormalTok{  )}

\NormalTok{df }\OtherTok{\textless{}{-}}\NormalTok{ loyalty\_clean }\SpecialCharTok{\%\textgreater{}\%}
  \FunctionTok{left\_join}\NormalTok{(customer\_activity, }\AttributeTok{by =} \StringTok{"Loyalty Number"}\NormalTok{) }\SpecialCharTok{\%\textgreater{}\%}
  \FunctionTok{mutate}\NormalTok{(}
    \AttributeTok{is\_promo\_2018 =} \StringTok{\textasciigrave{}}\AttributeTok{Enrollment Type}\StringTok{\textasciigrave{}} \SpecialCharTok{==} \StringTok{"2018 Promotion"}\NormalTok{,}
    \AttributeTok{churned =} \SpecialCharTok{!}\FunctionTok{is.na}\NormalTok{(}\StringTok{\textasciigrave{}}\AttributeTok{Cancellation Year}\StringTok{\textasciigrave{}}\NormalTok{)}
\NormalTok{  )}

\FunctionTok{glimpse}\NormalTok{(df)}
\end{Highlighting}
\end{Shaded}

\begin{verbatim}
## Rows: 16,737
## Columns: 26
## $ `Loyalty Number`      <dbl> 480934, 549612, 429460, 608370, 530508, 193662, ~
## $ Country               <chr> "Canada", "Canada", "Canada", "Canada", "Canada"~
## $ Province              <chr> "Ontario", "Alberta", "British Columbia", "Ontar~
## $ City                  <chr> "Toronto", "Edmonton", "Vancouver", "Toronto", "~
## $ `Postal Code`         <chr> "M2Z 4K1", "T3G 6Y6", "V6E 3D9", "P1W 1K4", "J8Y~
## $ Gender                <chr> "Female", "Male", "Male", "Male", "Male", "Male"~
## $ Education             <chr> "Bachelor", "College", "College", "College", "Ba~
## $ Salary                <dbl> 83236, NA, NA, NA, 103495, 51124, NA, 100159, 10~
## $ `Marital Status`      <chr> "Married", "Divorced", "Single", "Single", "Marr~
## $ `Loyalty Card`        <chr> "Star", "Star", "Star", "Star", "Star", "Star", ~
## $ CLV                   <dbl> 3839.14, 3839.61, 3839.75, 3839.75, 3842.79, 384~
## $ `Enrollment Type`     <chr> "Standard", "Standard", "Standard", "Standard", ~
## $ `Enrollment Year`     <dbl> 2016, 2016, 2014, 2013, 2014, 2012, 2014, 2016, ~
## $ `Enrollment Month`    <dbl> 2, 3, 7, 2, 10, 5, 6, 12, 5, 10, 3, 6, 1, 10, 11~
## $ `Cancellation Year`   <dbl> NA, NA, 2018, NA, NA, NA, NA, NA, NA, NA, NA, NA~
## $ `Cancellation Month`  <dbl> NA, NA, 1, NA, NA, NA, NA, NA, NA, NA, NA, NA, N~
## $ total_flights_sum     <dbl> 37, 58, 18, 35, 37, 0, 40, 41, 27, 44, 51, 31, 2~
## $ total_distance        <dbl> 54525, 77487, 24803, 48432, 55515, 0, 61007, 728~
## $ points_accumulated    <dbl> 54525, 77487, 24803, 48432, 55515, 0, 61007, 728~
## $ points_redeemed       <dbl> 1418, 1971, 374, 1291, 0, 0, 1583, 733, 1516, 16~
## $ dollar_cost_redeemed  <dbl> 256, 356, 67, 233, 0, 0, 285, 132, 274, 303, 90,~
## $ active_months         <int> 24, 24, 24, 24, 24, 4, 24, 24, 24, 24, 24, 24, 2~
## $ redemption_rate       <dbl> 0.026006419, 0.025436525, 0.015078821, 0.0266559~
## $ avg_flights_per_month <dbl> 1.5416667, 2.4166667, 0.7500000, 1.4583333, 1.54~
## $ is_promo_2018         <lgl> FALSE, FALSE, FALSE, FALSE, FALSE, FALSE, FALSE,~
## $ churned               <lgl> FALSE, FALSE, TRUE, FALSE, FALSE, FALSE, FALSE, ~
\end{verbatim}

Sedaj nimamo nobenih podvojenih vrednosti in negativnih plač, prav tako
pa smo v customer activity združili osebne podatke o strankah in
povzetek njihovih letalskih aktivnosti.

\section{Vizualizacija}\label{vizualizacija}

\section{(Manca) Razumeti dejavnike, ki vplivajo na vključenost v
program zvestobe, in analizirati učinkovitost promocijske
kampanje}\label{manca-razumeti-dejavnike-ki-vplivajo-na-vkljuux10denost-v-program-zvestobe-in-analizirati-uux10dinkovitost-promocijske-kampanje}

\subsection{Učinkovitost 2018
promocije}\label{uux10dinkovitost-2018-promocije}

\begin{Shaded}
\begin{Highlighting}[]
\DocumentationTok{\#\# {-}{-}{-}{-}{-}{-}{-}{-}{-}{-}{-}{-}{-}{-}{-}{-}{-}{-}{-}{-}{-}{-}{-}{-}{-}{-}{-}{-}{-}{-}{-}{-}{-}{-}{-}{-}{-}{-}{-}{-}{-}{-}{-}{-}{-}{-}{-}{-}{-}{-}{-}{-}{-}{-}{-}{-}{-}{-}{-}{-}{-}{-}{-}{-}{-}{-}{-}{-}{-}{-}{-}{-}{-}}
\DocumentationTok{\#\# 4. UČINKOVITOST 2018 PROMOCIJE (primerjava skupin)}
\DocumentationTok{\#\# {-}{-}{-}{-}{-}{-}{-}{-}{-}{-}{-}{-}{-}{-}{-}{-}{-}{-}{-}{-}{-}{-}{-}{-}{-}{-}{-}{-}{-}{-}{-}{-}{-}{-}{-}{-}{-}{-}{-}{-}{-}{-}{-}{-}{-}{-}{-}{-}{-}{-}{-}{-}{-}{-}{-}{-}{-}{-}{-}{-}{-}{-}{-}{-}{-}{-}{-}{-}{-}{-}{-}{-}{-}}

\CommentTok{\# Opisna statistika po skupinah}
\NormalTok{df }\SpecialCharTok{\%\textgreater{}\%}
  \FunctionTok{group\_by}\NormalTok{(}\StringTok{\textasciigrave{}}\AttributeTok{Enrollment Type}\StringTok{\textasciigrave{}}\NormalTok{) }\SpecialCharTok{\%\textgreater{}\%}
  \FunctionTok{summarise}\NormalTok{(}
    \AttributeTok{n =} \FunctionTok{n}\NormalTok{(),}
    \AttributeTok{povp\_flights   =} \FunctionTok{mean}\NormalTok{(total\_flights\_sum, }\AttributeTok{na.rm =} \ConstantTok{TRUE}\NormalTok{),}
    \AttributeTok{povp\_distance  =} \FunctionTok{mean}\NormalTok{(total\_distance, }\AttributeTok{na.rm =} \ConstantTok{TRUE}\NormalTok{),}
    \AttributeTok{povp\_CLV       =} \FunctionTok{mean}\NormalTok{(CLV, }\AttributeTok{na.rm =} \ConstantTok{TRUE}\NormalTok{),}
    \AttributeTok{povp\_accumulated =} \FunctionTok{mean}\NormalTok{(points\_accumulated, }\AttributeTok{na.rm =} \ConstantTok{TRUE}\NormalTok{),}
    \AttributeTok{povp\_redeemed     =} \FunctionTok{mean}\NormalTok{(points\_redeemed, }\AttributeTok{na.rm =} \ConstantTok{TRUE}\NormalTok{),}
    \AttributeTok{povp\_redemption\_rate =} \FunctionTok{mean}\NormalTok{(redemption\_rate, }\AttributeTok{na.rm =} \ConstantTok{TRUE}\NormalTok{),}
    \AttributeTok{povp\_dollar\_redeemed =} \FunctionTok{mean}\NormalTok{(dollar\_cost\_redeemed, }\AttributeTok{na.rm =} \ConstantTok{TRUE}\NormalTok{),}
    \AttributeTok{churn\_rate     =} \FunctionTok{mean}\NormalTok{(churned, }\AttributeTok{na.rm =} \ConstantTok{TRUE}\NormalTok{)}
\NormalTok{  )}
\end{Highlighting}
\end{Shaded}

\begin{verbatim}
## # A tibble: 2 x 10
##   `Enrollment Type`     n povp_flights povp_distance povp_CLV povp_accumulated
##   <chr>             <int>        <dbl>         <dbl>    <dbl>            <dbl>
## 1 2018 Promotion      971         45.8        68827.    8047.          103240.
## 2 Standard          15766         29.0        43437.    7985.           43437.
## # i 4 more variables: povp_redeemed <dbl>, povp_redemption_rate <dbl>,
## #   povp_dollar_redeemed <dbl>, churn_rate <dbl>
\end{verbatim}

\begin{Shaded}
\begin{Highlighting}[]
\CommentTok{\# Statistični test: CLV med Standard in 2018 Promotion (Wilcoxon, ker CLV ni nujno normalno porazdeljen)}
\FunctionTok{wilcox.test}\NormalTok{(CLV }\SpecialCharTok{\textasciitilde{}} \StringTok{\textasciigrave{}}\AttributeTok{Enrollment Type}\StringTok{\textasciigrave{}}\NormalTok{, }\AttributeTok{data =}\NormalTok{ df)}
\end{Highlighting}
\end{Shaded}

\begin{verbatim}
## 
##  Wilcoxon rank sum test with continuity correction
## 
## data:  CLV by Enrollment Type
## W = 7683302, p-value = 0.8432
## alternative hypothesis: true location shift is not equal to 0
\end{verbatim}

\begin{Shaded}
\begin{Highlighting}[]
\CommentTok{\# t{-}test za primerjavo povprečnega števila poletov}
\FunctionTok{t.test}\NormalTok{(total\_flights\_sum }\SpecialCharTok{\textasciitilde{}} \StringTok{\textasciigrave{}}\AttributeTok{Enrollment Type}\StringTok{\textasciigrave{}}\NormalTok{, }\AttributeTok{data =}\NormalTok{ df)}
\end{Highlighting}
\end{Shaded}

\begin{verbatim}
## 
##  Welch Two Sample t-test
## 
## data:  total_flights_sum by Enrollment Type
## t = 24.987, df = 1042.6, p-value < 0.00000000000000022
## alternative hypothesis: true difference in means between group 2018 Promotion and group Standard is not equal to 0
## 95 percent confidence interval:
##  15.54703 18.19698
## sample estimates:
## mean in group 2018 Promotion       mean in group Standard 
##                     45.84758                     28.97558
\end{verbatim}

\begin{Shaded}
\begin{Highlighting}[]
\CommentTok{\# Chi{-}kvadrat: ali je tip kartice povezan z enrollment type?}
\FunctionTok{chisq.test}\NormalTok{(}\FunctionTok{table}\NormalTok{(df}\SpecialCharTok{$}\StringTok{\textasciigrave{}}\AttributeTok{Enrollment Type}\StringTok{\textasciigrave{}}\NormalTok{, df}\SpecialCharTok{$}\StringTok{\textasciigrave{}}\AttributeTok{Loyalty Card}\StringTok{\textasciigrave{}}\NormalTok{))}
\end{Highlighting}
\end{Shaded}

\begin{verbatim}
## 
##  Pearson's Chi-squared test
## 
## data:  table(df$`Enrollment Type`, df$`Loyalty Card`)
## X-squared = 0.68438, df = 2, p-value = 0.7102
\end{verbatim}

\begin{Shaded}
\begin{Highlighting}[]
\CommentTok{\# Testi razlik med skupinama}
\FunctionTok{wilcox.test}\NormalTok{(points\_accumulated }\SpecialCharTok{\textasciitilde{}} \StringTok{\textasciigrave{}}\AttributeTok{Enrollment Type}\StringTok{\textasciigrave{}}\NormalTok{, }\AttributeTok{data =}\NormalTok{ df)}
\end{Highlighting}
\end{Shaded}

\begin{verbatim}
## 
##  Wilcoxon rank sum test with continuity correction
## 
## data:  points_accumulated by Enrollment Type
## W = 13176777, p-value < 0.00000000000000022
## alternative hypothesis: true location shift is not equal to 0
\end{verbatim}

\begin{Shaded}
\begin{Highlighting}[]
\FunctionTok{wilcox.test}\NormalTok{(points\_redeemed }\SpecialCharTok{\textasciitilde{}} \StringTok{\textasciigrave{}}\AttributeTok{Enrollment Type}\StringTok{\textasciigrave{}}\NormalTok{, }\AttributeTok{data =}\NormalTok{ df)}
\end{Highlighting}
\end{Shaded}

\begin{verbatim}
## 
##  Wilcoxon rank sum test with continuity correction
## 
## data:  points_redeemed by Enrollment Type
## W = 5434755, p-value < 0.00000000000000022
## alternative hypothesis: true location shift is not equal to 0
\end{verbatim}

\begin{Shaded}
\begin{Highlighting}[]
\FunctionTok{wilcox.test}\NormalTok{(redemption\_rate }\SpecialCharTok{\textasciitilde{}} \StringTok{\textasciigrave{}}\AttributeTok{Enrollment Type}\StringTok{\textasciigrave{}}\NormalTok{, }\AttributeTok{data =}\NormalTok{ df)}
\end{Highlighting}
\end{Shaded}

\begin{verbatim}
## 
##  Wilcoxon rank sum test with continuity correction
## 
## data:  redemption_rate by Enrollment Type
## W = 4248464, p-value < 0.00000000000000022
## alternative hypothesis: true location shift is not equal to 0
\end{verbatim}

Na prvi pogled nam opisna statistika pokaže, da je promocija iz leta
2018 privabila 971 novih članov, ki so v obstoju naročnine leteli v
povprečju kar 45-krat, medtem ko so normalne stranke letele le 29-krat v
povprečju. Zelo hitro je jasno, da je promocija pripeljala stranke, ki
so bolj pogosti uporabniki letalskih družb, kot standardno pridobljeni
kartičarji.

Testi za primerjavo skupin po načinu vpisa (standardno ali promocijsko)
kažejo naslednje rezultate. Stranke, ki so se vključile preko promocije,
so v opazovanjem obdobju v povprečju opravile skoraj 17 letov več od
standardno vpisanih strank. Ta razlika je statistično visoko značilna,
saj Welchov t-test vrne rezultate t = 24.987 in p-value \textless{}
2.2e-16. Povprečni vrednosti vidni v tabeli pa sta za promocijsko
skupino 45.8 in za standardno skupino 28.9.

Po drugi strani pa se CLV (customer lifetime value) med skupinama
praktično ne razlikuje, saj so rezultati Wilcoxonovega testa W = 7683302
in p-value = 0.8432. Promocija torej ni pritegnila strank z višjo
dolgoročno vrednostjo, temveč zgolj bolj aktivne letalce.

Na nadaljnih testih vidimo, da je zboranje in koriščenje točk visoko
značilno med skupinama, vendar pa nam samo ta rezultat ne pove, katera
skupina je bistveno boljša. Zato se obrnemo na tabelo s primerjavo, kjer
opazimo, da medtem ko promocijski člani zberejo bistveno več točk
zvestobe, so standardno včlanjeni člani tisti, ki jih več porabijo.
Redemption rate je v povprečju pri standardnih članih 1.7\%, medtem ko
je pri promocijskih le 0.4\%.

Zato je zanimivo, da je promocija privabila stranke, ki letijo veliko,
se pa ne ozirajo dosti na koriščenje ugodnosti, ki jih program nudi.
Potrebna bi bila nadaljna obravnava po plačilnem razredu in izobrazbi,
vendar sklep je, da stranke, ki veliko letijo nimajo finančnih zadržkov
in jim tako točke zvestobe niso pomemben del programa. Morda so jim bolj
važne druge ugodnosti, ki jih program zvestobe nudi, vendar teh podatkov
pa žal nimamo.

Tip kartice zvestobe (Star/Nova/Aurora) pa ni statistično povezan s
tipom vpisa, saj je Chi test vrnil rezultete X-squared = 0.68438 in
p-value = 0.7102. Promocija torej ni bila bolj privlačna za določen
segment kartic.

Pomembna opomba je, da ker gre za opazovalne podatke, višja aktivnost
promocijske skupine ne pomeni nujno, da je promocija vzročno povečala
angažiranost. Enako verjeten scenarij je samoselekcija, in sicer da so
se bolj pogosti letalci raje odločili za vključitev preko promocije. Za
bolj zanesljiv vzročni zaključek bi bilo treba primerjati angažiranost
istih strank pred in po vpisu.

Primerjaj v tabeli iz poglavja 4: pri Standard skupini je povp\_distance
= 43.437 in povp\_accumulated = 43.437 --- popolnoma identično. Pri 2018
Promotion skupini pa je povp\_distance = 68.827, povp\_accumulated =
103.240, kar je natanko 1,5-kratnik. To ni naključje --- nakazuje, da
imajo promocijske stranke 1,5× množitelj točk na prevoženi km,
standardne pa 1×. To je pomembna ugotovitev za tvojo nalogo
``učinkovitost promocije'': pomeni, da promocijska kampanja verjetno ni
bila zgolj marketinška (privabi nove člane), ampak je vključevala tudi
strukturno bonus stopnjo nabiranja točk kot del ponudbe. To delno
pojasni, zakaj promocijski člani nabirajo tako veliko več točk --- ne le
zaradi več letov, temveč tudi zaradi ugodnejšega razmerja točk/km.
Priporočam, da to eksplicitno preveriš in izpostaviš:

\begin{Shaded}
\begin{Highlighting}[]
\NormalTok{df }\SpecialCharTok{\%\textgreater{}\%}
  \FunctionTok{group\_by}\NormalTok{(}\StringTok{\textasciigrave{}}\AttributeTok{Enrollment Type}\StringTok{\textasciigrave{}}\NormalTok{) }\SpecialCharTok{\%\textgreater{}\%}
  \FunctionTok{summarise}\NormalTok{(}\AttributeTok{tocke\_na\_km =} \FunctionTok{sum}\NormalTok{(points\_accumulated, }\AttributeTok{na.rm =} \ConstantTok{TRUE}\NormalTok{) }\SpecialCharTok{/} \FunctionTok{sum}\NormalTok{(total\_distance, }\AttributeTok{na.rm =} \ConstantTok{TRUE}\NormalTok{))}
\end{Highlighting}
\end{Shaded}

\begin{verbatim}
## # A tibble: 2 x 2
##   `Enrollment Type` tocke_na_km
##   <chr>                   <dbl>
## 1 2018 Promotion            1.5
## 2 Standard                  1
\end{verbatim}

To bi lepo dopolnilo tvoj odstavek o tem, zakaj promocijski člani
``zberejo bistveno več točk, a jih manj unovčijo'' --- zdaj imaš razlago
tudi za zakaj zbirajo toliko več, ne le hipotezo o vedenju.

Manjkajo: - pogled tudi po izobrazbi in prihodku.

\subsection{Faktorji, ki vplivajo na
engagement}\label{faktorji-ki-vplivajo-na-engagement}

\begin{Shaded}
\begin{Highlighting}[]
\DocumentationTok{\#\# {-}{-}{-}{-}{-}{-}{-}{-}{-}{-}{-}{-}{-}{-}{-}{-}{-}{-}{-}{-}{-}{-}{-}{-}{-}{-}{-}{-}{-}{-}{-}{-}{-}{-}{-}{-}{-}{-}{-}{-}{-}{-}{-}{-}{-}{-}{-}{-}{-}{-}{-}{-}{-}{-}{-}{-}{-}{-}{-}{-}{-}{-}{-}{-}{-}{-}{-}{-}{-}{-}{-}{-}{-}}
\DocumentationTok{\#\# 5. FAKTORJI, KI VPLIVAJO NA ENGAGEMENT}
\DocumentationTok{\#\# {-}{-}{-}{-}{-}{-}{-}{-}{-}{-}{-}{-}{-}{-}{-}{-}{-}{-}{-}{-}{-}{-}{-}{-}{-}{-}{-}{-}{-}{-}{-}{-}{-}{-}{-}{-}{-}{-}{-}{-}{-}{-}{-}{-}{-}{-}{-}{-}{-}{-}{-}{-}{-}{-}{-}{-}{-}{-}{-}{-}{-}{-}{-}{-}{-}{-}{-}{-}{-}{-}{-}{-}{-}}

\CommentTok{\# Korelacija numeričnih spremenljivk z engagementom}
\NormalTok{df }\SpecialCharTok{\%\textgreater{}\%}
  \FunctionTok{select}\NormalTok{(Salary, CLV, total\_flights\_sum, total\_distance, redemption\_rate) }\SpecialCharTok{\%\textgreater{}\%}
  \FunctionTok{cor}\NormalTok{(}\AttributeTok{use =} \StringTok{"pairwise.complete.obs"}\NormalTok{)}
\end{Highlighting}
\end{Shaded}

\begin{verbatim}
##                         Salary          CLV total_flights_sum total_distance
## Salary             1.000000000 -0.021016986      -0.001653776   -0.002261964
## CLV               -0.021016986  1.000000000      -0.006916046   -0.005992790
## total_flights_sum -0.001653776 -0.006916046       1.000000000    0.967839881
## total_distance    -0.002261964 -0.005992790       0.967839881    1.000000000
## redemption_rate    0.009981686 -0.006658419       0.015660545   -0.004235114
##                   redemption_rate
## Salary                0.009981686
## CLV                  -0.006658419
## total_flights_sum     0.015660545
## total_distance       -0.004235114
## redemption_rate       1.000000000
\end{verbatim}

\begin{Shaded}
\begin{Highlighting}[]
\FunctionTok{library}\NormalTok{(reshape2)  }\CommentTok{\# za melt(), ali uporabi tidyr::pivot\_longer}
\end{Highlighting}
\end{Shaded}

\begin{verbatim}
## 
## Attaching package: 'reshape2'
\end{verbatim}

\begin{verbatim}
## The following objects are masked from 'package:data.table':
## 
##     dcast, melt
\end{verbatim}

\begin{verbatim}
## The following object is masked from 'package:tidyr':
## 
##     smiths
\end{verbatim}

\begin{Shaded}
\begin{Highlighting}[]
\NormalTok{cor\_matrix }\OtherTok{\textless{}{-}}\NormalTok{ df }\SpecialCharTok{\%\textgreater{}\%}
  \FunctionTok{select}\NormalTok{(Salary, CLV, total\_flights\_sum, total\_distance, redemption\_rate) }\SpecialCharTok{\%\textgreater{}\%}
  \FunctionTok{cor}\NormalTok{(}\AttributeTok{use =} \StringTok{"pairwise.complete.obs"}\NormalTok{)}

\NormalTok{cor\_long }\OtherTok{\textless{}{-}} \FunctionTok{melt}\NormalTok{(cor\_matrix)}

\FunctionTok{ggplot}\NormalTok{(cor\_long, }\FunctionTok{aes}\NormalTok{(}\AttributeTok{x =}\NormalTok{ Var1, }\AttributeTok{y =}\NormalTok{ Var2, }\AttributeTok{fill =}\NormalTok{ value)) }\SpecialCharTok{+}
  \FunctionTok{geom\_tile}\NormalTok{(}\AttributeTok{color =} \StringTok{"white"}\NormalTok{) }\SpecialCharTok{+}
  \FunctionTok{geom\_text}\NormalTok{(}\FunctionTok{aes}\NormalTok{(}\AttributeTok{label =} \FunctionTok{round}\NormalTok{(value, }\DecValTok{2}\NormalTok{)), }\AttributeTok{color =} \StringTok{"black"}\NormalTok{, }\AttributeTok{size =} \DecValTok{4}\NormalTok{) }\SpecialCharTok{+}
  \FunctionTok{scale\_fill\_gradient2}\NormalTok{(}\AttributeTok{low =} \StringTok{"steelblue"}\NormalTok{, }\AttributeTok{mid =} \StringTok{"white"}\NormalTok{, }\AttributeTok{high =} \StringTok{"firebrick"}\NormalTok{,}
                        \AttributeTok{midpoint =} \DecValTok{0}\NormalTok{, }\AttributeTok{limits =} \FunctionTok{c}\NormalTok{(}\SpecialCharTok{{-}}\DecValTok{1}\NormalTok{, }\DecValTok{1}\NormalTok{), }\AttributeTok{name =} \StringTok{"r"}\NormalTok{) }\SpecialCharTok{+}
  \FunctionTok{labs}\NormalTok{(}\AttributeTok{title =} \StringTok{"Korelacijska matrika: Salary, CLV in engagement metrike"}\NormalTok{,}
       \AttributeTok{x =} \ConstantTok{NULL}\NormalTok{, }\AttributeTok{y =} \ConstantTok{NULL}\NormalTok{) }\SpecialCharTok{+}
  \FunctionTok{theme\_minimal}\NormalTok{() }\SpecialCharTok{+}
  \FunctionTok{theme}\NormalTok{(}\AttributeTok{axis.text.x =} \FunctionTok{element\_text}\NormalTok{(}\AttributeTok{angle =} \DecValTok{45}\NormalTok{, }\AttributeTok{hjust =} \DecValTok{1}\NormalTok{))}
\end{Highlighting}
\end{Shaded}

\pandocbounded{\includegraphics[keepaspectratio]{Projekt---Airline-Loyalty-Program-Promotion-and-Beyond_files/Projekt---Airline-Loyalty-Program-Promotion-and-Beyond_files/figure-latex/unnamed-chunk-14-1.pdf}}

Na grafu korelacije vidimo, da ni absolutno nobene korelacije, ki bi
pojasnila aktivnost.

\begin{Shaded}
\begin{Highlighting}[]
\CommentTok{\# Linearni model: kaj napoveduje število poletov?}
\NormalTok{model }\OtherTok{\textless{}{-}} \FunctionTok{lm}\NormalTok{(total\_flights\_sum }\SpecialCharTok{\textasciitilde{}}\NormalTok{ Salary }\SpecialCharTok{+}\NormalTok{ Education }\SpecialCharTok{+} \StringTok{\textasciigrave{}}\AttributeTok{Marital Status}\StringTok{\textasciigrave{}} \SpecialCharTok{+}
\NormalTok{              Gender }\SpecialCharTok{+} \StringTok{\textasciigrave{}}\AttributeTok{Loyalty Card}\StringTok{\textasciigrave{}} \SpecialCharTok{+}\NormalTok{ is\_promo\_2018, }\AttributeTok{data =}\NormalTok{ df)}
\FunctionTok{summary}\NormalTok{(model)}
\end{Highlighting}
\end{Shaded}

\begin{verbatim}
## 
## Call:
## lm(formula = total_flights_sum ~ Salary + Education + `Marital Status` + 
##     Gender + `Loyalty Card` + is_promo_2018, data = df)
## 
## Residuals:
##     Min      1Q  Median      3Q     Max 
## -47.131 -11.438   3.879  11.925  58.121 
## 
## Coefficients:
##                                   Estimate   Std. Error t value
## (Intercept)                   27.842309870  0.670703444  41.512
## Salary                         0.000004200  0.000006268   0.670
## EducationDoctor               -0.216336591  0.911056343  -0.237
## EducationHigh School or Below  0.453265122  0.609818543   0.743
## EducationMaster                0.485052968  0.773201749   0.627
## `Marital Status`Married        0.377154059  0.399097535   0.945
## `Marital Status`Single         0.281474003  0.505460426   0.557
## GenderMale                     0.186506474  0.292140790   0.638
## `Loyalty Card`Nova             0.168137710  0.403374806   0.417
## `Loyalty Card`Star             0.519243046  0.387278518   1.341
## is_promo_2018TRUE             16.937215090  0.629640551  26.900
##                                          Pr(>|t|)    
## (Intercept)                   <0.0000000000000002 ***
## Salary                                      0.503    
## EducationDoctor                             0.812    
## EducationHigh School or Below               0.457    
## EducationMaster                             0.530    
## `Marital Status`Married                     0.345    
## `Marital Status`Single                      0.578    
## GenderMale                                  0.523    
## `Loyalty Card`Nova                          0.677    
## `Loyalty Card`Star                          0.180    
## is_promo_2018TRUE             <0.0000000000000002 ***
## ---
## Signif. codes:  0 '***' 0.001 '**' 0.01 '*' 0.05 '.' 0.1 ' ' 1
## 
## Residual standard error: 16.31 on 12467 degrees of freedom
##   (4259 observations deleted due to missingness)
## Multiple R-squared:  0.05511,    Adjusted R-squared:  0.05435 
## F-statistic: 72.71 on 10 and 12467 DF,  p-value: < 0.00000000000000022
\end{verbatim}

Korelacijska matrika pokaže, da so povezave med Salary, CLV in metrikami
angažiranosti (št. poletov, razdalja, redemption rate) praktično ničelne
(vse \textbar r\textbar{} \textless{} 0,03) -- plača in CLV torej sama
po sebi ne napovedujeta, koliko stranka leti. Pričakovano je le zelo
visoka korelacija med total\_flights\_sum in total\_distance (r = 0,97),
kar je logično, saj gre za dve meri istega pojava (več poletov = več
prepotovanih km), zato ju v modelu ne smemo uporabljati skupaj kot
neodvisni spremenljivki.

Linearni model (total\_flights\_sum \textasciitilde{} Salary + Education
+ Marital Status + Gender + Loyalty Card + is\_promo\_2018) potrjuje ta
vtis: edini statistično značilen napovednik je is\_promo\_2018 (+17,1
poleta, p \textless{} 0,001). Vse demografske spremenljivke -- plača,
izobrazba, zakonski stan, spol in tip kartice -- niso statistično
značilne napovednice angažiranosti. Model kot celota pojasni le 5,7 \%
variance (R² = 0,057), kar je nizko in kaže, da glavni dejavniki
angažiranosti niso zajeti v razpoložljivih demografskih spremenljivkah,
temveč verjetno v vedenjskih ali časovnih dejavnikih (npr. namen
potovanja, sezona, življenjski dogodki), ki jih v tem naboru podatkov
nimamo.

Sklep: demografija (kdo je stranka) pojasnjuje presenetljivo malo o tem,
kako pogosto stranka leti -- edini dejavnik z jasnim in močnim učinkom
je pripadnost 2018 promocijski kohorti.

\subsection{Vizualizacija}\label{vizualizacija-1}

\begin{Shaded}
\begin{Highlighting}[]
\DocumentationTok{\#\# {-}{-}{-}{-}{-}{-}{-}{-}{-}{-}{-}{-}{-}{-}{-}{-}{-}{-}{-}{-}{-}{-}{-}{-}{-}{-}{-}{-}{-}{-}{-}{-}{-}{-}{-}{-}{-}{-}{-}{-}{-}{-}{-}{-}{-}{-}{-}{-}{-}{-}{-}{-}{-}{-}{-}{-}{-}{-}{-}{-}{-}{-}{-}{-}{-}{-}{-}{-}{-}{-}{-}{-}{-}}
\DocumentationTok{\#\# 6. VIZUALIZACIJE}
\DocumentationTok{\#\# {-}{-}{-}{-}{-}{-}{-}{-}{-}{-}{-}{-}{-}{-}{-}{-}{-}{-}{-}{-}{-}{-}{-}{-}{-}{-}{-}{-}{-}{-}{-}{-}{-}{-}{-}{-}{-}{-}{-}{-}{-}{-}{-}{-}{-}{-}{-}{-}{-}{-}{-}{-}{-}{-}{-}{-}{-}{-}{-}{-}{-}{-}{-}{-}{-}{-}{-}{-}{-}{-}{-}{-}{-}}

\CommentTok{\# 6a. Distribucija CLV po enrollment tipu (boxplot {-} primerjava skupin, robusten na osamelce)}
\FunctionTok{ggplot}\NormalTok{(df, }\FunctionTok{aes}\NormalTok{(}\AttributeTok{x =} \StringTok{\textasciigrave{}}\AttributeTok{Enrollment Type}\StringTok{\textasciigrave{}}\NormalTok{, }\AttributeTok{y =}\NormalTok{ CLV, }\AttributeTok{fill =} \StringTok{\textasciigrave{}}\AttributeTok{Enrollment Type}\StringTok{\textasciigrave{}}\NormalTok{)) }\SpecialCharTok{+}
  \FunctionTok{geom\_boxplot}\NormalTok{(}\AttributeTok{outlier.alpha =} \FloatTok{0.3}\NormalTok{) }\SpecialCharTok{+}
  \FunctionTok{labs}\NormalTok{(}\AttributeTok{title =} \StringTok{"CLV glede na tip vpisa v program zvestobe"}\NormalTok{,}
       \AttributeTok{y =} \StringTok{"Customer Lifetime Value"}\NormalTok{, }\AttributeTok{x =} \ConstantTok{NULL}\NormalTok{) }\SpecialCharTok{+}
  \FunctionTok{theme\_minimal}\NormalTok{() }\SpecialCharTok{+}
  \FunctionTok{theme}\NormalTok{(}\AttributeTok{legend.position =} \StringTok{"none"}\NormalTok{)}
\end{Highlighting}
\end{Shaded}

\pandocbounded{\includegraphics[keepaspectratio]{Projekt---Airline-Loyalty-Program-Promotion-and-Beyond_files/Projekt---Airline-Loyalty-Program-Promotion-and-Beyond_files/figure-latex/unnamed-chunk-16-1.pdf}}

\begin{Shaded}
\begin{Highlighting}[]
\CommentTok{\# 6b. Časovna vrsta: povprečno število poletov na mesec po letih (učinek promocije v času)}
\NormalTok{flight\_clean }\SpecialCharTok{\%\textgreater{}\%}
  \FunctionTok{group\_by}\NormalTok{(Year, Month) }\SpecialCharTok{\%\textgreater{}\%}
  \FunctionTok{summarise}\NormalTok{(}\AttributeTok{povp\_flights =} \FunctionTok{mean}\NormalTok{(}\StringTok{\textasciigrave{}}\AttributeTok{Total Flights}\StringTok{\textasciigrave{}}\NormalTok{, }\AttributeTok{na.rm =} \ConstantTok{TRUE}\NormalTok{), }\AttributeTok{.groups =} \StringTok{"drop"}\NormalTok{) }\SpecialCharTok{\%\textgreater{}\%}
  \FunctionTok{mutate}\NormalTok{(}\AttributeTok{date =} \FunctionTok{make\_date}\NormalTok{(Year, Month, }\DecValTok{1}\NormalTok{)) }\SpecialCharTok{\%\textgreater{}\%}
  \FunctionTok{ggplot}\NormalTok{(}\FunctionTok{aes}\NormalTok{(}\AttributeTok{x =}\NormalTok{ date, }\AttributeTok{y =}\NormalTok{ povp\_flights)) }\SpecialCharTok{+}
  \FunctionTok{geom\_line}\NormalTok{(}\AttributeTok{color =} \StringTok{"steelblue"}\NormalTok{) }\SpecialCharTok{+}
  \FunctionTok{geom\_vline}\NormalTok{(}\AttributeTok{xintercept =} \FunctionTok{as.numeric}\NormalTok{(}\FunctionTok{make\_date}\NormalTok{(}\DecValTok{2018}\NormalTok{, }\DecValTok{1}\NormalTok{, }\DecValTok{1}\NormalTok{)),}
             \AttributeTok{linetype =} \StringTok{"dashed"}\NormalTok{, }\AttributeTok{color =} \StringTok{"red"}\NormalTok{) }\SpecialCharTok{+}
  \FunctionTok{labs}\NormalTok{(}\AttributeTok{title =} \StringTok{"Povprečno število poletov skozi čas"}\NormalTok{,}
       \AttributeTok{subtitle =} \StringTok{"Rdeča črta = začetek 2018 promocije"}\NormalTok{,}
       \AttributeTok{x =} \ConstantTok{NULL}\NormalTok{, }\AttributeTok{y =} \StringTok{"Povp. Total Flights"}\NormalTok{) }\SpecialCharTok{+}
  \FunctionTok{theme\_minimal}\NormalTok{()}
\end{Highlighting}
\end{Shaded}

\begin{verbatim}
## Warning in scale_x_date(): A <numeric> value was passed to a Date scale.
## i The value was converted to a <Date> object.
\end{verbatim}

\pandocbounded{\includegraphics[keepaspectratio]{Projekt---Airline-Loyalty-Program-Promotion-and-Beyond_files/Projekt---Airline-Loyalty-Program-Promotion-and-Beyond_files/figure-latex/unnamed-chunk-16-2.pdf}}

\begin{Shaded}
\begin{Highlighting}[]
\CommentTok{\# 6c. Redemption rate po tipu kartice in enrollment tipu}
\FunctionTok{ggplot}\NormalTok{(df, }\FunctionTok{aes}\NormalTok{(}\AttributeTok{x =} \StringTok{\textasciigrave{}}\AttributeTok{Loyalty Card}\StringTok{\textasciigrave{}}\NormalTok{, }\AttributeTok{y =}\NormalTok{ redemption\_rate, }\AttributeTok{fill =} \StringTok{\textasciigrave{}}\AttributeTok{Enrollment Type}\StringTok{\textasciigrave{}}\NormalTok{)) }\SpecialCharTok{+}
  \FunctionTok{geom\_boxplot}\NormalTok{() }\SpecialCharTok{+}
  \FunctionTok{labs}\NormalTok{(}\AttributeTok{title =} \StringTok{"Stopnja unovčitve točk po tipu kartice in vpisa"}\NormalTok{,}
       \AttributeTok{y =} \StringTok{"Redemption rate"}\NormalTok{, }\AttributeTok{x =} \StringTok{"Loyalty Card"}\NormalTok{) }\SpecialCharTok{+}
  \FunctionTok{theme\_minimal}\NormalTok{()}
\end{Highlighting}
\end{Shaded}

\begin{verbatim}
## Warning: Removed 1 row containing non-finite outside the scale range
## (`stat_boxplot()`).
\end{verbatim}

\pandocbounded{\includegraphics[keepaspectratio]{Projekt---Airline-Loyalty-Program-Promotion-and-Beyond_files/Projekt---Airline-Loyalty-Program-Promotion-and-Beyond_files/figure-latex/unnamed-chunk-16-3.pdf}}

\begin{Shaded}
\begin{Highlighting}[]
\CommentTok{\# 6d. Churn rate po enrollment tipu (bar chart {-} deleži)}
\NormalTok{df }\SpecialCharTok{\%\textgreater{}\%}
  \FunctionTok{group\_by}\NormalTok{(}\StringTok{\textasciigrave{}}\AttributeTok{Enrollment Type}\StringTok{\textasciigrave{}}\NormalTok{) }\SpecialCharTok{\%\textgreater{}\%}
  \FunctionTok{summarise}\NormalTok{(}\AttributeTok{churn\_rate =} \FunctionTok{mean}\NormalTok{(churned)) }\SpecialCharTok{\%\textgreater{}\%}
  \FunctionTok{ggplot}\NormalTok{(}\FunctionTok{aes}\NormalTok{(}\AttributeTok{x =} \StringTok{\textasciigrave{}}\AttributeTok{Enrollment Type}\StringTok{\textasciigrave{}}\NormalTok{, }\AttributeTok{y =}\NormalTok{ churn\_rate, }\AttributeTok{fill =} \StringTok{\textasciigrave{}}\AttributeTok{Enrollment Type}\StringTok{\textasciigrave{}}\NormalTok{)) }\SpecialCharTok{+}
  \FunctionTok{geom\_col}\NormalTok{() }\SpecialCharTok{+}
  \FunctionTok{scale\_y\_continuous}\NormalTok{(}\AttributeTok{labels =}\NormalTok{ scales}\SpecialCharTok{::}\NormalTok{percent) }\SpecialCharTok{+}
  \FunctionTok{labs}\NormalTok{(}\AttributeTok{title =} \StringTok{"Delež odpovedi članstva po tipu vpisa"}\NormalTok{, }\AttributeTok{y =} \StringTok{"Churn rate"}\NormalTok{, }\AttributeTok{x =} \ConstantTok{NULL}\NormalTok{) }\SpecialCharTok{+}
  \FunctionTok{theme\_minimal}\NormalTok{() }\SpecialCharTok{+}
  \FunctionTok{theme}\NormalTok{(}\AttributeTok{legend.position =} \StringTok{"none"}\NormalTok{)}
\end{Highlighting}
\end{Shaded}

\pandocbounded{\includegraphics[keepaspectratio]{Projekt---Airline-Loyalty-Program-Promotion-and-Beyond_files/Projekt---Airline-Loyalty-Program-Promotion-and-Beyond_files/figure-latex/unnamed-chunk-16-4.pdf}}

\section{(Tristan) Identificirajte segmente strank z visoko vrednostjo
in predlagajte možne ukrepe za
optimizacijo}\label{tristan-identificirajte-segmente-strank-z-visoko-vrednostjo-in-predlagajte-moux17ene-ukrepe-za-optimizacijo}

\section{(Patricija) Napovedovanje odliva strank oziroma odpovedi ter
predlaganje možnih ukrepov za njihovo
zadržanje}\label{patricija-napovedovanje-odliva-strank-oziroma-odpovedi-ter-predlaganje-moux17enih-ukrepov-za-njihovo-zadrux17eanje}

\end{document}
